\documentclass[11pt,a4paper,openany]{book}

% ============================================================
% Packages
% ============================================================
\usepackage[utf8]{inputenc}
\usepackage[T1]{fontenc}
\usepackage{lmodern}
\usepackage[margin=2.5cm]{geometry}
\usepackage{amsmath,amssymb,amsthm,mathtools}
\usepackage{bm}
\usepackage{graphicx}
\usepackage{xcolor}
\usepackage{hyperref}
\usepackage{enumitem}
\usepackage{fancyhdr}
\usepackage{tikz}
\usepackage{microtype}
\hypersetup{colorlinks=true, linkcolor=blue!60!black, citecolor=blue!60!black,
            urlcolor=blue!60!black}

% ============================================================
% Theorem environments
% ============================================================
\newtheoremstyle{plainserif}{10pt}{10pt}{\itshape}{}%
   {\bfseries}{.}{0.5em}{}
\newtheoremstyle{remarkstyle}{10pt}{10pt}{}{}%
   {\itshape}{.}{0.5em}{}

\theoremstyle{plainserif}
\newtheorem{theorem}{Theorem}[chapter]
\newtheorem{proposition}[theorem]{Proposition}
\newtheorem{lemma}[theorem]{Lemma}
\newtheorem{corollary}[theorem]{Corollary}

\theoremstyle{definition}
\newtheorem{definition}[theorem]{Definition}
\newtheorem{observation}[theorem]{Observation}
\newtheorem{conjecture}[theorem]{Conjecture}

\theoremstyle{remarkstyle}
\newtheorem*{remark}{Remark}
\newtheorem*{interpretation}{Interpretation}

% Side-note box for ``why'' explanations
\definecolor{whybox}{RGB}{245,245,230}
\newenvironment{whybox}{%
  \begin{center}\begin{minipage}{0.95\linewidth}\small
  \textbf{\textit{Why this step?}}\quad\ignorespaces}{%
  \end{minipage}\end{center}}

% ============================================================
% Macros
% ============================================================
\newcommand{\R}{\mathbb{R}}
\newcommand{\N}{\mathbb{N}}
\newcommand{\E}{\mathbb{E}}
\newcommand{\V}{\operatorname{Var}}
\newcommand{\Cov}{\operatorname{Cov}}
\newcommand{\Normal}{\mathcal{N}}
\newcommand{\Lap}{\mathrm{Lap}}
\newcommand{\btilde}{\tilde{b}}
\newcommand{\hatfn}[1]{\widehat{#1}}
\newcommand{\mills}{R}
\newcommand{\tr}{\operatorname{tr}}
\DeclareMathOperator{\Hess}{Hess}
\newcommand{\ind}{\mathbf{1}}

% Header style
\pagestyle{fancy}
\fancyhf{}
\fancyhead[LE]{\small \leftmark}
\fancyhead[RO]{\small \rightmark}
\fancyfoot[C]{\small \thepage}
\renewcommand{\chaptermark}[1]{\markboth{\small\textsc{Chapter \thechapter.\ #1}}{}}
\renewcommand{\sectionmark}[1]{\markright{\small\textsc{\thesection.\ #1}}{}}

% ============================================================
% Title page
% ============================================================
\title{\bfseries Mathematical Compendium for the
Laplace AR(1) Diffusion Benchmark\\[4pt]
{\large\itshape Full derivations, proofs, and interpretations accompanying the
main memo}}
\author{Giovanni Mantovani \and Marco Lomel\'e \\[6pt]
\small Supervisor: Prof.\ Marc M\'ezard \quad Tutor: J\'er\^ome Garnier-Brun\\[2pt]
\small Bocconi University \textendash\ MSc Data Science \& AI}
\date{\today}

\begin{document}
\frontmatter
\maketitle

% ============================================================
% Preface
% ============================================================
\chapter*{Preface}
\addcontentsline{toc}{chapter}{Preface}

This compendium accompanies the main memo
\emph{``Exact Benchmark for the Laplace AR(1) Diffusion''} and provides the
full mathematical workings behind every statement made there.
The main memo is designed to be read quickly; it states the results, shows
the key formulas, and presents the numerical validation. It assumes the
reader trusts the derivations.
This document does not make that assumption. Every theorem is proved,
every algebraic step is written out, every change of variables is justified,
and every interpretation is anchored in the preceding algebra.

The style is deliberately pedagogical. When a step may look surprising,
we pause to explain why it works --- why a particular substitution is made,
why a particular coordinate is natural, why a particular limit gives the
behaviour we claim. This is meant to be a document one could hand to a
collaborator (or to oneself six months from now) and recover the full
mathematical story without having to reconstruct any step.

We chose a book format rather than a long article because the material
divides naturally into six self-contained parts:

\begin{description}[leftmargin=1.5em,style=nextline]
\item[Part I --- Setting and foundations.]
  The generative model, the noising channel, and the structural constants
  that the whole analysis rests on.
\item[Part II --- The $K=1$ structure theorems.]
  Factorisation, residual density, joint score, curvature field. The
  core of the exact benchmark.
\item[Part III --- Fourier representation.]
  The characteristic function, the rotation to innovation-frequency
  coordinates, the linear ODE satisfied by the residual density.
\item[Part IV --- Scalar observables.]
  The Mills-ratio closed form for the centre curvature, Gaussian
  benchmark, kurtosis, FWHM.
\item[Part V --- The $K=2$ case.]
  The exact attempt, the obstruction to $1$D factorisation, and the
  exact four-term representation.
\item[Part VI --- Interpretation.]
  What Gaussianity buys us, what Markovianity buys us, and where the
  two structural roles can be disentangled.
\end{description}

Throughout, we keep the same notation as the original derivation
note~\cite{laplace_note} (reproduced in Table~\ref{tab:constants} below)
so that this document reads as a careful expansion of that note rather
than a parallel account.

\paragraph{How to read this document.}
The main theorems are numbered and cross-referenced with the memo.
Proofs are written out in full; each non-trivial algebraic step that we
believe a reader could stumble on is boxed as a ``why this step'' note.
An equation labelled (e.g.) \textbf{Eq.~\eqref{eq:h-closed}} in the memo
has the same reference here, so that one can flip between documents.

\tableofcontents

% ============================================================
\mainmatter
% ============================================================

\part{Setting and foundations}

% ============================================================
\chapter{The generative model and the question we are asking}
\label{ch:setting}
% ============================================================

\section{The clean chain}

We study an autoregressive chain of order one driven by Laplace
innovations. Fix parameters $\alpha \in \R$, $b > 0$, $\sigma_0 > 0$,
$\mu_0 \in \R$, and define
\begin{equation}
A_0 \sim \Normal(\mu_0, \sigma_0^2), \qquad
A_{k+1} = \alpha A_k + \varepsilon_{k+1},\quad
\varepsilon_k \stackrel{\mathrm{iid}}{\sim} \Lap(0, b),
\label{eq:clean-chain}
\end{equation}
where $\Lap(0,b)$ has density $(2b)^{-1}\exp(-|u|/b)$. The chain
$(A_k)_{k\geq 0}$ is a first-order Markov chain: conditional on $A_k$, the
future $A_{k+1}, A_{k+2}, \ldots$ is independent of the past
$A_0, \ldots, A_{k-1}$. The innovation $\varepsilon_k$ is the part of
$A_k$ that is new -- not predictable from $A_{k-1}$.

We refer to $(A_k)$ as the \emph{clean chain}. The density of
$(A_0, \ldots, A_{K-1})$ factorises:
\begin{equation}
p_0(a_0, \ldots, a_{K-1}) = \phi_{\mu_0,\sigma_0^2}(a_0)\,
\prod_{k=1}^{K-1} \frac{1}{2b}\exp\!\left(-\frac{|a_k - \alpha a_{k-1}|}{b}\right),
\label{eq:clean-density}
\end{equation}
where $\phi_{m,\sigma^2}$ denotes the Gaussian pdf with mean $m$ and
variance $\sigma^2$. This factorisation reflects the clean Markov property:
the joint density is a product over bonds, each bond depending only on two
consecutive variables.

\begin{whybox}
  Equation~\eqref{eq:clean-density} is what makes the AR(1) chain a
  Markov random field with a tridiagonal coupling structure. If we were
  to work in Fourier space, each bond would contribute a rational factor
  of the form $1/(1 + b^2 \ell^2)$ to the characteristic function along
  its own innovation-frequency axis, a structure we will exploit
  heavily in Part III.
\end{whybox}

\section{The OU corruption channel}

Generative diffusion proceeds by smoothly noising the clean data through
the Ornstein--Uhlenbeck (OU) channel:
\begin{equation}
X_k(t) = \beta\, A_k + \sqrt{\Delta}\, Z_k,\qquad
\beta := e^{-t},\qquad \Delta := 1 - e^{-2t},\qquad
Z_k \stackrel{\mathrm{iid}}{\sim} \Normal(0, 1),
\label{eq:OU}
\end{equation}
with $(Z_0, \ldots, Z_{K-1})$ independent of the chain
$(A_0, \ldots, A_{K-1})$. For each fixed $t$, the transformation
$A_k \mapsto X_k(t)$ is affine-Gaussian and coordinatewise: each
$X_k(t)$ depends only on $A_k$ and an independent Gaussian. In
particular, conditional on $A = (A_0, \ldots, A_{K-1})$, the
coordinates $X_0(t), \ldots, X_{K-1}(t)$ are independent Gaussians.

This is the coordinate-local variant of the standard variance-preserving
diffusion kernel; an alternative view starts from the SDE
$\mathrm{d}X_k = -X_k\,\mathrm{d}t + \sqrt{2}\,\mathrm{d}B_k$ with
initial condition $X_k(0)=A_k$, whose marginal distribution at time $t$
is exactly \eqref{eq:OU}.

\begin{whybox}
Three properties of the OU channel will be used over and over again:
(i) at $t=0$ we have $\beta=1$, $\Delta=0$ and $X=A$ identically;
(ii) as $t\to\infty$, $\beta \to 0$, $\Delta\to 1$ and $X \to \Normal(0,I)$
regardless of the data;
(iii) the channel is \emph{conditionally independent across coordinates}
given the clean state.
Property (iii) is the reason Markov structure survives noising, as we
will prove in Section~\ref{sec:markov-survives}.
\end{whybox}

\section{The central question}
\label{sec:central-question}

Given the noisy distribution at time $t$,
$p_t(x_0, x_1, \ldots, x_{K-1})$, we study the \emph{joint score}
\begin{equation}
S(x, t) := \nabla_x \log p_t(x_0, \ldots, x_{K-1}) \in \R^K,
\qquad x = (x_0, \ldots, x_{K-1}),
\label{eq:joint-score}
\end{equation}
and the \emph{curvature field}
\begin{equation}
H_t(x) := -\nabla_x^2 \log p_t(x_0, \ldots, x_{K-1}) \in \R^{K\times K}.
\label{eq:curvature}
\end{equation}
These are different from the per-frame marginal scores
$\nabla_{x_k} \log p_t(x_k)$: the latter discard correlations between
different time steps and therefore miss the whole point of diffusion over
dynamic objects. Our question is:

\begin{quote}
\emph{How does the joint score $S(x, t)$ reflect the Markov structure of
the clean chain, and what happens to that structure as $t$ varies?}
\end{quote}

In the Gaussian AR(1) benchmark (studied in~\cite{structure_sigma_t}), the
answer is known exactly: $p_t$ is Gaussian, $S$ is linear, and
$H_t = \Sigma_t^{-1}$ is a constant precision matrix that interpolates
between tridiagonal (at $t=0$, reflecting the Markov structure) and a
full matrix (for large $t$). The present compendium works out the
answer in the simplest non-Gaussian case -- the Laplace chain -- where
Gaussianity and Markovianity play separable structural roles. The
resulting exact benchmark then serves as a fulcrum for the thesis: any
general-purpose theory of joint scores over dynamic objects must recover
this benchmark in its appropriate limits.

\begin{remark}[Joint vs.\ marginal scores]
Throughout, ``score'' means joint score as in~\eqref{eq:joint-score}.
Whenever we need the per-frame marginal score
$\nabla_{x_k} \log p_t(x_k)$ we will write ``marginal score.''
\end{remark}

% ============================================================
\chapter{Why Markov structure survives noising}
\label{ch:markov-survives}
% ============================================================

Before diving into explicit computations, we establish a structural fact
that underpins all of Part II: the noisy law $p_t$ is itself a Markov
random field over time, with exactly the same sparsity graph as the clean
law. This is not automatic -- it requires the conditional independence
property of the OU channel. The proof is short and instructive.

\label{sec:markov-survives}
\begin{proposition}[Markov property survives OU noising]
\label{prop:markov-survives}
Under~\eqref{eq:clean-chain}--\eqref{eq:OU}, the noisy sequence
$(X_0, \ldots, X_{K-1})$ is Markov of order one: for every $k$,
\begin{equation}
p_t(x_k \mid x_0, \ldots, x_{k-1}) = p_t(x_k \mid x_{k-1}).
\label{eq:noisy-markov}
\end{equation}
Equivalently, $X_{k+1} \perp (X_0, \ldots, X_{k-1}) \mid X_k$.
\end{proposition}

\begin{proof}
We will use two facts repeatedly:
\begin{itemize}[leftmargin=2em, topsep=3pt]
\item[(i)] \textbf{Clean Markov property.}
$A_{k+1} \perp (A_0,\ldots,A_{k-1}) \mid A_k$.
\item[(ii)] \textbf{Channel independence.}
Given the entire clean chain $A$, the noises $Z_0,\ldots,Z_{K-1}$
are i.i.d.\ $\Normal(0,1)$ and independent of $A$, so
$(X_0, \ldots, X_{K-1}) \mid A$ has independent coordinates and
$X_k \mid A = A_k$ depends only on $A_k$.
\end{itemize}
We compute $p_t(x_k \mid x_0, \ldots, x_{k-1})$ by conditioning on the clean
chain:
\begin{align}
p_t(x_k \mid x_0, \ldots, x_{k-1})
 &= \int p_t(x_k \mid A, x_0, \ldots, x_{k-1}) \,
        p(A \mid x_0, \ldots, x_{k-1}) \,\mathrm{d}A
 \label{eq:markov-pf-1}\\
 &= \int p_t(x_k \mid A_k)\,
        p(A \mid x_0, \ldots, x_{k-1}) \,\mathrm{d}A,
 \label{eq:markov-pf-2}
\end{align}
where the second equality uses channel independence (ii): given $A$,
$X_k$ depends only on $A_k$ and is independent of
$X_0,\ldots,X_{k-1}$. After integrating out $A_{k+1},\ldots,A_{K-1}$
(which are irrelevant) and the components $A_0, \ldots, A_{k-1}$ (which
do not appear in $p_t(x_k \mid A_k)$), we are left with
\begin{equation}
p_t(x_k \mid x_0, \ldots, x_{k-1}) =
\int p_t(x_k \mid A_k)\, p(A_k \mid x_0, \ldots, x_{k-1}) \,\mathrm{d}A_k.
\label{eq:markov-pf-3}
\end{equation}
It remains to show that
$p(A_k \mid x_0,\ldots,x_{k-1}) = p(A_k \mid x_{k-1})$.
Condition on $A_{k-1}$ and use the clean Markov property~(i):
$A_k \perp (A_0,\ldots,A_{k-2}) \mid A_{k-1}$.
Channel independence gives
$A_k \perp (X_0,\ldots,X_{k-2}) \mid A_{k-1}$,
and furthermore
$(A_k, A_{k-1}) \perp (X_0,\ldots,X_{k-2}) \mid X_{k-1}$
because of the memoryless Markov structure of
$A_0 \to X_0$, $A_1 \to X_1$, etc.
Integrating out $A_{k-1}$ then yields
$p(A_k \mid x_0,\ldots,x_{k-1}) = p(A_k \mid x_{k-1})$, so
\eqref{eq:markov-pf-3} reduces to
\begin{equation*}
p_t(x_k \mid x_0, \ldots, x_{k-1}) =
\int p_t(x_k \mid A_k)\, p(A_k \mid x_{k-1}) \,\mathrm{d}A_k
= p_t(x_k \mid x_{k-1}),
\end{equation*}
which is~\eqref{eq:noisy-markov}.
\end{proof}

\begin{interpretation}
The noisy chain is genuinely Markov: every noisy transition
$p_t(x_{k+1} \mid x_k)$ encodes the full conditional dependence structure
at time $t$. The deep reason is that the OU channel is
\emph{coordinate-local}: it never creates new coupling between coordinates
at fixed $t$. Any coupling in $(X_0, \ldots, X_{K-1})$ must be inherited
from the clean chain.
\end{interpretation}

\begin{corollary}[Factorisation of the noisy joint density]
\label{cor:noisy-factor}
We have
\begin{equation}
p_t(x_0, \ldots, x_{K-1}) = p_t(x_0)\, \prod_{k=1}^{K-1} p_t(x_k \mid x_{k-1}).
\label{eq:noisy-factor}
\end{equation}
\end{corollary}

\begin{remark}[Score structure]
Because the joint log-density is a sum of two-point terms,
$\log p_t = \log p_t(x_0) + \sum_{k} \log p_t(x_k \mid x_{k-1})$,
the Hessian $H_t = -\nabla^2 \log p_t$ is a sum of block matrices each of
size $2\times 2$ supported on adjacent indices. This is why the joint
score and curvature field have a tridiagonal block structure at $K\geq 2$
even before we compute anything explicitly.
\end{remark}

% ============================================================
\chapter{Structural constants for $K=1$}
\label{ch:constants}
% ============================================================

We now specialise to $K=1$, i.e.\ we work with the pair $(X_0, X_1)$ noised
from $(A_0, A_1)$. This already contains the non-trivial physics: a Gaussian
prior, a Laplace innovation, and a coordinate-local Gaussian channel.
The goal of this chapter is to compute, from first principles, all the
structural constants that appear in the exact formulas -- the
reparameterisations $v$, $s^2$, $\tau^2$, $\rho$, $\nu$, $\btilde$, $a$.
Every later derivation rests on these, so we are careful to show where
each quantity comes from.

\section{The marginal $p_t(x_0)$}

Because $X_0 = \beta A_0 + \sqrt{\Delta}Z_0$ with
$A_0 \sim \Normal(\mu_0, \sigma_0^2)$ and $Z_0 \sim \Normal(0,1)$
independent,
\begin{equation}
X_0 \sim \Normal(\beta\mu_0, v), \qquad
v := \beta^2 \sigma_0^2 + \Delta,
\label{eq:x0-marginal}
\end{equation}
and the density is
\begin{equation}
p_t(x_0) = \frac{1}{\sqrt{2\pi v}}\exp\!\left(-\frac{(x_0-\beta\mu_0)^2}{2v}\right),
\qquad
\partial_{x_0}\log p_t(x_0) = -\frac{x_0 - \beta\mu_0}{v}.
\label{eq:x0-score}
\end{equation}

\begin{whybox}
The quantity $v$ interpolates between $\sigma_0^2$ (at $t=0$, where
$\beta=1,\Delta=0$) and $1$ (as $t\to\infty$, where $\beta\to 0,\Delta\to 1$).
It is the total variance of $X_0$, combining the prior variance
(scaled down by $\beta^2$) and the injected OU noise variance $\Delta$.
\end{whybox}

\section{The posterior $A_0 \mid X_0 = x_0$}

The joint $(A_0, X_0)$ is bivariate Gaussian, so the posterior
$A_0 \mid X_0 = x_0$ is again Gaussian. The standard conditioning
formulas are
\begin{align}
m(x_0) &:= \E[A_0 \mid X_0 = x_0]
  = \mu_0 + \frac{\Cov(A_0, X_0)}{\V(X_0)}(x_0 - \E X_0), \\
s^2 &:= \V(A_0 \mid X_0 = x_0)
  = \V(A_0) - \frac{\Cov(A_0, X_0)^2}{\V(X_0)}.
\end{align}
We have
$\Cov(A_0, X_0) = \Cov(A_0, \beta A_0 + \sqrt{\Delta}Z_0) = \beta\sigma_0^2$,
so
\begin{align}
m(x_0) &= \mu_0 + \frac{\beta \sigma_0^2}{v}(x_0 - \beta\mu_0)
       = \frac{\Delta \mu_0 + \beta\sigma_0^2 x_0}{v},
\label{eq:posterior-mean}\\
s^2 &= \sigma_0^2 - \frac{\beta^2 \sigma_0^4}{v} = \frac{\sigma_0^2 \Delta}{v}.
\label{eq:posterior-var}
\end{align}

\begin{whybox}
Equation~\eqref{eq:posterior-mean} can be rewritten as
$m(x_0) = \mu_0 + \rho'\, (x_0 - \beta\mu_0)$
with $\rho' = \beta\sigma_0^2/v$; this is the familiar
``shrinkage factor'' of Gaussian regression.
Equation~\eqref{eq:posterior-var} says the posterior variance
factorises into the prior variance $\sigma_0^2$ times the ratio
$\Delta/v$, which quantifies \emph{how much the noise veils the signal}:
it is zero at $t=0$ (we see $A_0$ perfectly), and it approaches
$\sigma_0^2$ as $t \to \infty$ (we see nothing).
\end{whybox}

\section{The residual random variable}
\label{sec:residual-setup}

We now derive the key identity
\begin{equation}
X_1 \mid X_0 = x_0 \stackrel{d}{=} \nu + \rho x_0 + R_t,
\qquad
R_t = \beta\varepsilon + G_t,\quad G_t \sim \Normal(0, \tau^2),
\label{eq:residual}
\end{equation}
where $R_t$ is independent of $X_0$ and of the particular $x_0$. The
constants $\nu, \rho, \tau^2$ will be pinned down by direct computation.

Start from
\[
X_1 = \beta A_1 + \sqrt{\Delta}Z_1 = \beta(\alpha A_0 + \varepsilon) + \sqrt{\Delta}Z_1
    = \beta\alpha A_0 + \beta\varepsilon + \sqrt{\Delta}Z_1.
\]
Decompose $A_0 = m(x_0) + \xi$ where $\xi := A_0 - m(x_0)$. Conditional on
$X_0 = x_0$, $\xi \sim \Normal(0, s^2)$ and, importantly, $\xi$ is
independent of $(\varepsilon, Z_1)$: indeed, $\xi$ is a function of
$(A_0, X_0) = (A_0, \beta A_0 + \sqrt{\Delta}Z_0)$, hence a linear combination
of $A_0$ and $Z_0$, both of which are independent of $\varepsilon$ and $Z_1$.
Therefore
\begin{equation}
X_1 \mid X_0 = x_0 \stackrel{d}{=} \beta\alpha m(x_0) + \bigl(\beta\alpha \xi + \beta\varepsilon
                    + \sqrt{\Delta}Z_1\bigr).
\label{eq:residual-pf-1}
\end{equation}
The deterministic part simplifies using~\eqref{eq:posterior-mean}:
\begin{align}
\beta\alpha m(x_0)
 &= \beta\alpha\cdot \frac{\Delta\mu_0 + \beta\sigma_0^2 x_0}{v}
  = \underbrace{\frac{\alpha\beta\Delta\mu_0}{v}}_{=:\, \nu}
    + \underbrace{\frac{\alpha\beta^2\sigma_0^2}{v}}_{=:\, \rho}\, x_0,
\label{eq:nu-rho}
\end{align}
which \emph{defines} the parameters $\nu$ and $\rho$.

The random part of~\eqref{eq:residual-pf-1} is
$R_t := \beta\alpha\xi + \beta\varepsilon + \sqrt{\Delta}Z_1$.
Split it as
\[
R_t = \beta\varepsilon + G_t, \qquad
G_t := \beta\alpha\xi + \sqrt{\Delta}Z_1.
\]
The variable $G_t$ is a linear combination of two independent Gaussians,
hence Gaussian with mean zero and variance
\begin{equation}
\tau^2 := \V(G_t) = \beta^2\alpha^2 s^2 + \Delta
       = \Delta + \alpha^2\beta^2\frac{\sigma_0^2\Delta}{v}
       = \Delta\left(1 + \frac{\alpha^2\beta^2\sigma_0^2}{v}\right).
\label{eq:tau2-derivation}
\end{equation}
Both $\varepsilon$ and $G_t$ are independent of each other: $\varepsilon$
is independent of $(\xi, Z_1)$ by construction and $G_t$ is built from
$(\xi, Z_1)$. Finally, the whole of $R_t$ is independent of $X_0$:
conditional on $x_0$, both $\varepsilon$ (by independence) and $G_t$ (with
mean $0$ and variance not depending on $x_0$) are independent of the
conditioning.

Collecting everything, we have proved~\eqref{eq:residual} with
\begin{equation}
\boxed{\;\rho = \dfrac{\alpha\beta^2\sigma_0^2}{v},\qquad
       \nu  = \dfrac{\alpha\beta\Delta\mu_0}{v},\qquad
       \tau^2 = \Delta + \alpha^2\beta^2 s^2,\qquad
       s^2 = \dfrac{\sigma_0^2\Delta}{v}.\;}
\label{eq:core-constants}
\end{equation}

We also define $\btilde := \beta b$ -- the effective Laplace scale of the
\emph{rescaled} innovation $\beta\varepsilon$, since $\beta\varepsilon \sim \Lap(0,\btilde)$.
The dimensionless quantity
\begin{equation}
a := \tau / \btilde
\label{eq:a-def}
\end{equation}
will appear repeatedly and plays the role of an inverse
non-Gaussianity parameter: large $a$ means the Gaussian part of $R_t$
dominates and $R_t$ is nearly Gaussian; small $a$ means the Laplace
part dominates and the cusp at $r=0$ is strongly felt.

\section{An auxiliary constant: the leakage $c$}

A slightly different parameterisation will be useful in Chapters~\ref{ch:fourier}
and~\ref{ch:interpretation}. Define the \emph{innovation coordinate}
\begin{equation}
y := x_1 - \alpha x_0,
\end{equation}
which is the noisy analogue of the clean innovation
$\varepsilon = A_1 - \alpha A_0$. Direct substitution gives
\begin{align}
Y := X_1 - \alpha X_0
 &= \beta(A_1 - \alpha A_0) + \sqrt{\Delta}(Z_1 - \alpha Z_0)
 = \beta\varepsilon + \sqrt{\Delta}(Z_1 - \alpha Z_0).
\label{eq:innovation-coord}
\end{align}
This variable contains no $A_0$ contribution at all -- its dependence on
$A$ is entirely through the innovation $\varepsilon$. But $Y$ is not
independent of $X_0$ at finite $t$, because both contain the noise $Z_0$.
To see how $y$ and the residual $r$ relate, compute
\begin{align}
r &= x_1 - \nu - \rho x_0 = y + \alpha x_0 - \nu - \rho x_0
   = y + (\alpha - \rho)\,x_0 - \nu.
\end{align}
Using~\eqref{eq:core-constants},
\begin{equation}
\alpha - \rho = \alpha - \frac{\alpha\beta^2\sigma_0^2}{v}
              = \alpha\cdot \frac{v - \beta^2\sigma_0^2}{v}
              = \alpha\cdot \frac{\Delta}{v} =: c,
\label{eq:c-def}
\end{equation}
so
\begin{equation}
r = y + c\,(x_0 - \beta\mu_0).
\label{eq:r-in-terms-of-y}
\end{equation}
The constant $c = \alpha\Delta/v$ measures how much the shared OU noise
``leaks'' $x_0$-dependence into the innovation channel. At $t=0$,
$\Delta = 0 \Rightarrow c = 0$, and the residual and innovation coincide;
at general $t$, a linear shift $c(x_0 - \beta\mu_0)$ separates them.

\begin{table}[h]
\centering\small
\begin{tabular}{l l l l}
\hline
Symbol & Definition & Meaning & Limit $t\to 0$ / $t\to\infty$\\ \hline
$\beta$ & $e^{-t}$ & OU shrinkage & $1$ / $0$\\
$\Delta$ & $1-\beta^2$ & OU injected variance & $0$ / $1$\\
$v$ & $\beta^2\sigma_0^2 + \Delta$ & $\V(X_0)$ & $\sigma_0^2$ / $1$\\
$s^2$ & $\sigma_0^2\Delta/v$ & $\V(A_0\mid X_0)$ & $0$ / $\sigma_0^2$\\
$\rho$ & $\alpha\beta^2\sigma_0^2/v$ & residual slope in $x_0$ & $\alpha$ / $0$\\
$\nu$ & $\alpha\beta\Delta\mu_0/v$ & residual offset & $0$ / $0$\\
$c$ & $\alpha\Delta/v$ & innovation leakage & $0$ / $\alpha$\\
$\btilde$ & $\beta b$ & rescaled Laplace scale & $b$ / $0$\\
$\tau^2$ & $\Delta + \alpha^2\beta^2 s^2$ & residual Gaussian variance & $0$ / $1$\\
$a$ & $\tau/\btilde$ & Mills argument & $0$ / $\infty$\\
\hline
\end{tabular}
\caption{Structural constants and their limits. All derived from
first principles in Chapter~\ref{ch:constants}.}
\label{tab:constants}
\end{table}

\begin{remark}[Check: the $t\to 0$ limit]
At $t=0$: $\beta=1$, $\Delta=0$, $v = \sigma_0^2$, $s^2 = 0$,
$\tau = 0$, $\rho = \alpha$, $\nu = 0$, $\btilde = b$, $a = 0$.
The residual becomes $R_0 = \varepsilon$, and~\eqref{eq:residual}
reduces to $X_1 = \alpha X_0 + \varepsilon$, i.e.\ the clean chain.
Good.
\end{remark}

\begin{remark}[Check: the $t\to\infty$ limit]
At $t\to\infty$: $\beta\to 0$, $\Delta\to 1$, $v\to 1$, $\tau^2\to 1$,
$\rho\to 0$, $\btilde\to 0$, $a\to\infty$. The residual
$R_t = \beta\varepsilon + G_t$ becomes a standard normal, and the chain
collapses to $X \sim \Normal(0, I)$, the diffusion prior. Good.
\end{remark}

% ============================================================
\part{The $K{=}1$ structure theorems}
% ============================================================

\chapter{Factorisation theorem}
\label{ch:fact}

The factorisation theorem states that the noisy joint density $p_t(x_0,x_1)$
decomposes exactly as a product of a Gaussian site density and a
one-dimensional non-Gaussian bond. Conceptually: the only non-Gaussian
content of $(X_0, X_1)$ lives in the one-dimensional residual coordinate
$r = x_1 - \nu - \rho x_0$; everything perpendicular to this direction
is Gaussian.

\begin{theorem}[Factorisation, $K{=}1$]
\label{thm:fact}
Under the model~\eqref{eq:clean-chain}--\eqref{eq:OU},
\begin{equation}
p_t(x_0, x_1) = g_t(x_0)\, h_t(r),
\qquad
r = x_1 - \nu - \rho x_0,
\label{eq:fact-main}
\end{equation}
where $g_t(x_0) = \frac{1}{\sqrt{2\pi v}}\exp\!\bigl(-(x_0-\beta\mu_0)^2/(2v)\bigr)$
is the Gaussian marginal of $X_0$, and $h_t$ is the density of the
Normal-Laplace random variable
\begin{equation}
R_t = \beta\varepsilon + G_t,\qquad
\varepsilon \sim \Lap(0,b),\;\;
G_t \sim \Normal(0,\tau^2),\;\;
\varepsilon \perp G_t.
\label{eq:Rt}
\end{equation}
\end{theorem}

\begin{proof}
By the chain rule of probability, $p_t(x_0, x_1) = p_t(x_0) p_t(x_1\mid x_0)$.
We have already shown $p_t(x_0) = g_t(x_0)$ in~\eqref{eq:x0-score}. From
Section~\ref{sec:residual-setup},
\[
X_1 \mid X_0 = x_0 \stackrel{d}{=} \nu + \rho x_0 + R_t,
\]
and $R_t$ is independent of $x_0$ (indeed, its distribution does not
depend on $x_0$ at all). By a standard change-of-variables argument, the
conditional density is
\begin{equation}
p_t(x_1 \mid x_0) = h_t(x_1 - \nu - \rho x_0) = h_t(r).
\end{equation}
Multiplying through gives~\eqref{eq:fact-main}.
\end{proof}

\begin{interpretation}[What this result is, and what it is not]
Theorem~\ref{thm:fact} is \emph{not} the clean factorisation of a Markov
chain -- that would read $p_t(x_0,x_1) = p_t(x_0) p_t(x_1\mid x_0)$, which
is always true and carries no information about the innovation structure.
What~\eqref{eq:fact-main} says is that the \emph{conditional law}
$p_t(x_1 \mid x_0)$ is a \emph{translate} of the same universal 1D density
$h_t$, at an $x_0$-dependent location $\nu + \rho x_0$. This is a
non-trivial consequence of the AR(1)+OU structure: both the deterministic
shift $\nu + \rho x_0$ and the shape $h_t$ are explicit.
\end{interpretation}

\begin{remark}
From this point on, every derivative of $\log p_t(x_0, x_1)$ with respect
to $x_0$ or $x_1$ can be computed via the chain rule through the single
variable $r$. This is why the subsequent theorems all read as ``take
$\psi_t := (\log h_t)'$, and here are the scores/Hessians.''
\end{remark}

% ============================================================
\chapter{The Normal-Laplace residual density}
\label{ch:convolution}
% ============================================================

The residual density $h_t$ is the convolution of a Laplace and a Gaussian
density. We compute it in closed form. The result is classical
(it is the Normal--Laplace density used in financial modelling), but it
is important to see the algebra because the \emph{form} of the answer --
two exponential terms, each multiplied by a Gaussian CDF -- is what makes
all subsequent derivations (score, curvature, Mills' ratio) tractable.

\section{Setup}

Since $\beta\varepsilon \sim \Lap(0, \btilde)$ with $\btilde = \beta b$, the density
of $R_t = \beta\varepsilon + G_t$ with $G_t \sim \Normal(0,\tau^2)$ independent is
\begin{equation}
h_t(r) = \int_{-\infty}^{\infty}
   \frac{1}{2\btilde}\, e^{-|u|/\btilde}\cdot
   \frac{1}{\sqrt{2\pi\tau^2}}\, e^{-(r-u)^2/(2\tau^2)} \,\mathrm{d}u.
\label{eq:h-convolution}
\end{equation}
The absolute value $|u|$ prevents a direct Gaussian integration, so we split
the real line at $u = 0$:
\begin{align}
h_t(r)
 &= \frac{1}{2\btilde\sqrt{2\pi\tau^2}}
   \left[ I_+(r) + I_-(r) \right],
\label{eq:h-split}
\end{align}
with
\begin{equation}
I_+(r) := \int_0^\infty\! \exp\!\left( -\frac{u}{\btilde} - \frac{(r-u)^2}{2\tau^2}\right)\!\mathrm{d}u,
\qquad
I_-(r) := \int_{-\infty}^0\!\exp\!\left( \frac{u}{\btilde} - \frac{(r-u)^2}{2\tau^2}\right)\!\mathrm{d}u.
\label{eq:I-pm-def}
\end{equation}
Each of $I_\pm$ is a half-line Gaussian integral with a linear term --
exactly the form that reduces to a complementary Gaussian CDF after
completing the square.

\section{Computation of $I_+(r)$}
\label{sec:I-plus}

Expand the exponent of the $I_+$ integrand:
\begin{align}
-\frac{u}{\btilde} - \frac{(r-u)^2}{2\tau^2}
&= -\frac{u}{\btilde} - \frac{u^2 - 2ru + r^2}{2\tau^2}
\label{eq:Iplus-expand}\\
&= -\frac{u^2}{2\tau^2} + \left(\frac{r}{\tau^2} - \frac{1}{\btilde}\right) u
   - \frac{r^2}{2\tau^2}.
\label{eq:Iplus-rearranged}
\end{align}
We now complete the square in $u$. Write
\begin{equation}
\frac{r}{\tau^2} - \frac{1}{\btilde} = \frac{r - \tau^2/\btilde}{\tau^2},
\end{equation}
so that the linear coefficient in $u$ equals $(r - \tau^2/\btilde)/\tau^2$.
The quadratic-plus-linear combination becomes
\begin{align}
-\frac{u^2}{2\tau^2} + \frac{r-\tau^2/\btilde}{\tau^2}\, u
&= -\frac{1}{2\tau^2}\left[ u^2 - 2(r - \tau^2/\btilde)\, u \right]
\label{eq:Iplus-sq-1}\\
&= -\frac{1}{2\tau^2}\left[ (u - r + \tau^2/\btilde)^2 - (r - \tau^2/\btilde)^2 \right].
\label{eq:Iplus-sq-2}
\end{align}

\begin{whybox}
Completing the square in $u$ converts ``Gaussian times exponential'' into
``Gaussian centred at a shifted point, times a deterministic factor.''
The shifted centre $u_* = r - \tau^2/\btilde$ tells us where the
integrand is concentrated; the deterministic factor $(r-\tau^2/\btilde)^2 - r^2$
absorbs the parts that do not depend on $u$.
\end{whybox}

Combining the completion-of-squares with the $-r^2/(2\tau^2)$ term from
\eqref{eq:Iplus-rearranged} gives
\begin{align}
-\frac{u}{\btilde} - \frac{(r-u)^2}{2\tau^2}
 &= -\frac{(u - r + \tau^2/\btilde)^2}{2\tau^2}
   + \frac{(r-\tau^2/\btilde)^2 - r^2}{2\tau^2}
\label{eq:Iplus-almost}\\
 &= -\frac{(u - r + \tau^2/\btilde)^2}{2\tau^2}
   - \frac{r}{\btilde} + \frac{\tau^2}{2\btilde^2}.
\label{eq:Iplus-final-exponent}
\end{align}
The last step uses
$\frac{(r-\tau^2/\btilde)^2 - r^2}{2\tau^2}
= \frac{-2r\tau^2/\btilde + \tau^4/\btilde^2}{2\tau^2}
= -\frac{r}{\btilde} + \frac{\tau^2}{2\btilde^2}$.

Pulling the deterministic factor out of the integral,
\begin{equation}
I_+(r) = e^{-r/\btilde + \tau^2/(2\btilde^2)}
   \int_0^\infty \exp\!\left(-\frac{(u - r + \tau^2/\btilde)^2}{2\tau^2}\right)\mathrm{d}u.
\label{eq:Iplus-after-square}
\end{equation}
Change variable: $z = (u - r + \tau^2/\btilde)/\tau$, so $\mathrm{d}u = \tau\,\mathrm{d}z$
and the lower limit $u = 0$ becomes $z = -r/\tau + \tau/\btilde$.
Using $a = \tau/\btilde$,
\begin{equation}
I_+(r) = e^{-r/\btilde + a^2/2}\, \tau\, \int_{-r/\tau + a}^\infty
         e^{-z^2/2}\,\mathrm{d}z
       = e^{-r/\btilde + a^2/2}\, \tau\sqrt{2\pi}\,\Phi\!\left(\frac{r}{\tau} - a\right).
\label{eq:Iplus-closed}
\end{equation}
Here $\Phi(x) := \int_{-\infty}^x \phi(z)\,\mathrm{d}z$ is the standard normal
CDF, and we used $\int_c^\infty e^{-z^2/2}\,\mathrm{d}z = \sqrt{2\pi}\,\Phi(-c)$
together with $\Phi(-c) = \Phi(c)\bigl|_{c\to -c}$.

\begin{remark}[Check the exponent constants]
The prefactor $\exp(-r/\btilde + a^2/2)$ has a natural interpretation:
$-r/\btilde$ is the Laplace exponential decay in the tail $r\to+\infty$,
and $a^2/2 = \tau^2/(2\btilde^2)$ is a finite correction coming from the
Gaussian smoothing. As $\tau\to 0$ the correction vanishes and we recover
pure Laplace behaviour.
\end{remark}

\section{Computation of $I_-(r)$}

By the symmetry $u \mapsto -u$,
\begin{align}
I_-(r) &= \int_{-\infty}^0 \exp\!\left(\frac{u}{\btilde} - \frac{(r-u)^2}{2\tau^2}\right)\mathrm{d}u
        = \int_0^\infty \exp\!\left(-\frac{u}{\btilde} - \frac{(r+u)^2}{2\tau^2}\right)\mathrm{d}u.
\end{align}
This is the same integral as $I_+$ with $r \to -r$. Therefore, by
\eqref{eq:Iplus-closed},
\begin{equation}
I_-(r) = e^{r/\btilde + a^2/2}\,\tau\sqrt{2\pi}\,\Phi\!\left(-\frac{r}{\tau} - a\right).
\label{eq:Iminus-closed}
\end{equation}

\section{The closed form of $h_t$}

Substituting~\eqref{eq:Iplus-closed} and~\eqref{eq:Iminus-closed} into
\eqref{eq:h-split} and simplifying the constants:
\begin{align}
h_t(r) &= \frac{1}{2\btilde\sqrt{2\pi\tau^2}}\cdot
   e^{a^2/2}\,\tau\sqrt{2\pi}\,
   \left[ e^{-r/\btilde}\Phi\!\left(\frac{r}{\tau}-a\right)
        + e^{r/\btilde}\Phi\!\left(-\frac{r}{\tau}-a\right) \right] \nonumber\\
&= \boxed{\;\frac{e^{a^2/2}}{2\btilde}
   \left[ e^{-r/\btilde}\,\Phi\!\left(\frac{r}{\tau}-a\right)
        + e^{r/\btilde}\,\Phi\!\left(-\frac{r}{\tau}-a\right) \right].\;}
\label{eq:h-closed}
\end{align}
This is the Normal--Laplace density; it is the key formula of Part II.

\begin{interpretation}
The two exponentials capture the two Laplace tails (decay to the right
and to the left), each damped by a Gaussian-CDF factor that smooths out
the cusp at $r = 0$. In the limit $\tau \to 0$ (so $a\to 0$), each
$\Phi(\cdot - a)$ becomes a Heaviside step that selects the correct half
of the Laplace, and we recover the clean Laplace density
$e^{-|r|/\btilde}/(2\btilde)$ as expected.
\end{interpretation}

\begin{remark}[Sanity check: $h_t$ integrates to 1]
As $r\to\pm\infty$, one of the two exponentials dominates and the other
vanishes; the Gaussian CDF in the dominant branch approaches $1$,
recovering pure Laplace decay. Integrating $e^{-|r|/\btilde}/(2\btilde)$
over $\R$ gives $1$, which is the leading contribution. The finite-$\tau$
corrections cancel; a direct computation via the convolution definition
confirms $\int h_t(r)\,\mathrm{d}r = 1$.
\end{remark}

\begin{definition}[$L_\pm$ and the compact form]
\label{def:Lpm}
We introduce the notation
\begin{equation}
L_+(r) := e^{-r/\btilde}\, \Phi\!\left(\frac{r}{\tau}-a\right),\qquad
L_-(r) := e^{r/\btilde}\, \Phi\!\left(-\frac{r}{\tau}-a\right),
\label{eq:Lpm-def}
\end{equation}
so that
\begin{equation}
h_t(r) = \frac{e^{a^2/2}}{2\btilde}\bigl[L_+(r) + L_-(r)\bigr].
\label{eq:h-compact}
\end{equation}
\end{definition}

% ============================================================
\chapter{Score theorem}
\label{ch:score}
% ============================================================

The next building block is the closed form of the joint score. We compute
the residual score $\psi_t = (\log h_t)'$ first; the full joint score
then follows by the chain rule through $r$.

\section{A useful identity}

Before differentiating, we record an identity that makes the computation
clean.

\begin{lemma}[Gaussian-derivative cancellation]
\label{lem:gaussian-id}
\begin{equation}
e^{-r/\btilde}\,\phi\!\left(\frac{r}{\tau}-a\right)
= e^{r/\btilde}\,\phi\!\left(\frac{r}{\tau}+a\right)
= e^{-a^2/2}\,\phi\!\left(\frac{r}{\tau}\right).
\label{eq:gauss-id}
\end{equation}
\end{lemma}

\begin{proof}
Using $\phi(x) = (2\pi)^{-1/2}e^{-x^2/2}$ and
$(r/\tau - a)^2 = (r/\tau)^2 - 2ar/\tau + a^2$, we get
$\phi(r/\tau - a) = \phi(r/\tau)\, e^{ar/\tau - a^2/2}$.
Since $a/\tau = 1/\btilde$,
\begin{equation}
e^{-r/\btilde}\,\phi(r/\tau - a) = e^{-r/\btilde}\,\phi(r/\tau)\,e^{r/\btilde - a^2/2}
  = e^{-a^2/2}\,\phi(r/\tau).
\end{equation}
The other equality is by the same argument with $r\to -r$.
\end{proof}

\begin{whybox}
This identity is at the heart of why $h_t'$ simplifies. When we
differentiate $L_+ + L_-$, each piece produces two terms: one from the
exponential part, one from the Gaussian CDF (whose derivative is a
Gaussian pdf). Lemma~\ref{lem:gaussian-id} says that the two
\emph{pdf} pieces are equal and cancel when summed with opposite signs.
The surviving terms come only from the exponentials.
\end{whybox}

\section{Derivative of $h_t$}

Differentiate $L_\pm$ using the product rule. Recall that
$\Phi'(x) = \phi(x)$, so $\frac{\mathrm{d}}{\mathrm{d}r}\Phi(r/\tau \pm a) = (1/\tau)\phi(r/\tau \pm a)$.
\begin{align}
L_+'(r) &= -\frac{1}{\btilde} e^{-r/\btilde}\,\Phi\!\left(\frac{r}{\tau}-a\right)
          + \frac{1}{\tau} e^{-r/\btilde}\,\phi\!\left(\frac{r}{\tau}-a\right),
\label{eq:Lplus-der}\\
L_-'(r) &= +\frac{1}{\btilde} e^{r/\btilde}\,\Phi\!\left(-\frac{r}{\tau}-a\right)
          - \frac{1}{\tau} e^{r/\btilde}\,\phi\!\left(\frac{r}{\tau}+a\right),
\label{eq:Lminus-der}
\end{align}
where we used $\phi(-r/\tau - a) = \phi(r/\tau + a)$ by symmetry of $\phi$
and a sign-flip under the inner derivative.
Now sum:
\begin{align}
L_+'(r) + L_-'(r)
 &= -\frac{1}{\btilde} L_+(r) + \frac{1}{\btilde} L_-(r)
    + \frac{1}{\tau}\Bigl[ e^{-r/\btilde}\phi(r/\tau - a)
       - e^{r/\btilde}\phi(r/\tau + a) \Bigr].
\end{align}
By Lemma~\ref{lem:gaussian-id}, the bracket in the last term is exactly
zero:
\begin{equation}
e^{-r/\btilde}\phi\!\left(\frac{r}{\tau} - a\right)
- e^{r/\btilde}\phi\!\left(\frac{r}{\tau} + a\right)
= e^{-a^2/2}\phi\!\left(\frac{r}{\tau}\right) - e^{-a^2/2}\phi\!\left(\frac{r}{\tau}\right) = 0.
\label{eq:pdf-cancel}
\end{equation}
Therefore
\begin{equation}
L_+'(r) + L_-'(r) = \frac{1}{\btilde}\bigl[L_-(r) - L_+(r)\bigr],
\label{eq:Lplus-plus-Lminus-deriv}
\end{equation}
and differentiating \eqref{eq:h-compact},
\begin{equation}
h_t'(r) = \frac{e^{a^2/2}}{2\btilde^2}\bigl[L_-(r) - L_+(r)\bigr].
\label{eq:hprime-closed}
\end{equation}

\section{The residual score}

Dividing $h_t'$ by $h_t$ and cancelling the $e^{a^2/2}/2$ factors:
\begin{equation}
\psi_t(r) := \frac{h_t'(r)}{h_t(r)} =
\frac{1}{\btilde}\,\frac{L_-(r) - L_+(r)}{L_+(r) + L_-(r)}.
\label{eq:psi-closed}
\end{equation}

\begin{theorem}[Joint score on the residual ray]
\label{thm:score}
The joint score $S(x,t) = \nabla_x \log p_t(x_0, x_1)$ of the $K=1$ model
equals
\begin{equation}
S(x,t) = \begin{pmatrix}
-\dfrac{x_0 - \beta\mu_0}{v} - \rho\,\psi_t(r) \\[4pt]
\psi_t(r)
\end{pmatrix},
\qquad r = x_1 - \nu - \rho x_0,
\label{eq:S-full}
\end{equation}
with $\psi_t$ given by~\eqref{eq:psi-closed}. In particular,
$\psi_t(r) \to \mp 1/\btilde$ as $r\to\pm\infty$.
\end{theorem}

\begin{proof}
From the factorisation $\log p_t(x_0,x_1) = \log g_t(x_0) + \log h_t(r)$,
\begin{align}
\partial_{x_1} \log p_t &= \partial_{x_1} \log h_t(r) = (\log h_t)'(r)\cdot \partial_{x_1} r = \psi_t(r),\\
\partial_{x_0} \log p_t &= (\log g_t)'(x_0) + (\log h_t)'(r)\cdot \partial_{x_0} r
                        = -\frac{x_0-\beta\mu_0}{v} + \psi_t(r)\cdot(-\rho),
\end{align}
where we used $\partial_{x_1} r = 1$ and $\partial_{x_0} r = -\rho$.
Stacking gives~\eqref{eq:S-full}.
\end{proof}

\subsection*{Tail behaviour of $\psi_t$}

We prove the saturation claim rigorously. Focus on $r\to+\infty$ (the other
side is identical by symmetry of $h_t$). Write the ratio in~\eqref{eq:psi-closed} as
\begin{equation}
\frac{L_-(r) - L_+(r)}{L_+(r) + L_-(r)}
 = \frac{L_-(r)/L_+(r) - 1}{L_-(r)/L_+(r) + 1}
 = \tanh\!\left(\tfrac{1}{2}\log\bigl(L_-(r)/L_+(r)\bigr)\right),
\label{eq:psi-tanh}
\end{equation}
where the second equality uses $(x-1)/(x+1) = \tanh((\log x)/2)$ for $x > 0$.
We have $L_-(r)/L_+(r) = e^{2r/\btilde}\,\Phi(-r/\tau - a)/\Phi(r/\tau - a)$.
As $r\to+\infty$: $\Phi(r/\tau - a) \to 1$ and $\Phi(-r/\tau - a) \to 0$,
but the factor $e^{2r/\btilde}$ grows so fast that we need to be careful.
Using Mills' ratio $\Phi(-x) \sim \phi(x)/x$ for $x\to +\infty$,
\begin{equation}
\Phi(-r/\tau - a) \sim \frac{\phi(r/\tau + a)}{r/\tau + a}\to 0 \text{ super-exponentially in }r,
\end{equation}
so $\log L_-(r) \sim r/\btilde - (r/\tau + a)^2/2 \sim -r^2/(2\tau^2)$ for large $r$,
while $\log L_+(r) \sim -r/\btilde$. Therefore
\begin{equation}
\log(L_-(r)/L_+(r)) \sim -\frac{r^2}{2\tau^2} + \frac{r}{\btilde} \to -\infty \text{ as }r\to\infty.
\end{equation}
Plugging into~\eqref{eq:psi-tanh}, $\tanh((-\infty)/2) = -1$, so
$\psi_t(r)\cdot \btilde \to -1$, i.e.\ $\psi_t(r) \to -1/\btilde$.
By $r\mapsto -r$ symmetry, $\psi_t(r)\to +1/\btilde$ as $r\to-\infty$. $\Box$

\begin{interpretation}[Why the score saturates]
For a Gaussian, the tails are so thin that escaping them costs quadratic
energy, and the restoring force -- the score -- grows linearly. For a Laplace,
the tails are only exponential; the cost of escape grows linearly, and the
restoring force is bounded. Here, even with a Gaussian contribution $G_t$
added to $\beta\varepsilon$, the Laplace tail dominates at large $|r|$ and the
score saturates at $\pm 1/\btilde = \pm 1/(\beta b)$. Saturation is the
signature of sub-Gaussian tails; any diffusion model that tries to
represent a Laplace-like innovation with a Gaussian score will be
\emph{locally} wrong in the tails by a factor growing in $|r|$.
\end{interpretation}

% ============================================================
\chapter{Curvature theorem}
\label{ch:curv}
% ============================================================

The residual score $\psi_t(r)$ depends only on $r$. The Hessian $H_t(x)$
is obtained by differentiating the score and thus inherits this
dependence: all of its entries are proportional to a single scalar
function $\kappa_t(r) := -\psi_t'(r)$. We show this explicitly and then
note the strong low-rank consequence.

\section{Componentwise computation}

From Theorem~\ref{thm:score}, we have
$S_0(x) = -(x_0-\beta\mu_0)/v - \rho\psi_t(r)$
and $S_1(x) = \psi_t(r)$. The Hessian is
\begin{equation}
H_t(x) = -\nabla_x^2 \log p_t = -\begin{pmatrix} \partial_{x_0} S_0 & \partial_{x_1} S_0 \\
                                                   \partial_{x_0} S_1 & \partial_{x_1} S_1 \end{pmatrix}.
\end{equation}
Using $\partial_{x_0} r = -\rho$ and $\partial_{x_1} r = 1$, and writing
$\psi_t'(r) = -\kappa_t(r)$,
\begin{align}
\partial_{x_1} S_1 &= \psi_t'(r) = -\kappa_t(r), \\
\partial_{x_0} S_1 &= \psi_t'(r)\cdot(-\rho) = \rho\kappa_t(r), \\
\partial_{x_1} S_0 &= -\rho\psi_t'(r) = \rho\kappa_t(r), \\
\partial_{x_0} S_0 &= -\frac{1}{v} - \rho\psi_t'(r)\cdot(-\rho) = -\frac{1}{v} - \rho^2 \kappa_t(r)
   \quad \text{(wait, with the sign of }\psi_t').
\end{align}
Let us redo the last one carefully:
$S_0 = -(x_0-\beta\mu_0)/v - \rho\psi_t(r)$, so
$\partial_{x_0} S_0 = -1/v - \rho \psi_t'(r)\cdot \partial_{x_0} r
                    = -1/v - \rho \psi_t'(r)\cdot(-\rho)
                    = -1/v + \rho^2 \psi_t'(r)
                    = -1/v - \rho^2 \kappa_t(r)$.
Hence
\begin{equation}
H_t(x) = -\begin{pmatrix} -1/v - \rho^2\kappa_t(r) & \rho\kappa_t(r) \\
                          \rho\kappa_t(r) & -\kappa_t(r) \end{pmatrix}
       = \begin{pmatrix} 1/v + \rho^2\kappa_t(r) & -\rho\kappa_t(r) \\
                          -\rho\kappa_t(r) & \kappa_t(r) \end{pmatrix}.
\label{eq:H-explicit}
\end{equation}

\begin{theorem}[Rank-1 structure of the curvature, $K{=}1$]
\label{thm:curv}
\begin{equation}
H_t(x) = \underbrace{\begin{pmatrix}1/v & 0\\0 & 0\end{pmatrix}}_{=:\, Q_0}
       \,+\, \kappa_t(r)\, M_\rho,\qquad
M_\rho := \begin{pmatrix}\rho^2 & -\rho\\ -\rho & 1\end{pmatrix}.
\label{eq:H-ray}
\end{equation}
The matrix $M_\rho$ has rank one, with $M_\rho = nn^\top$ where
$n = (\rho, -1)^\top$. Its eigenvalues are $\{0, 1+\rho^2\}$, with
eigenvectors $\propto (1,\rho)^\top$ (kernel) and $n$ respectively.
Consequently, for fixed $t$, the curvature field $H_t(x)$ traces a
one-parameter affine ray in the space of $2\times 2$ symmetric matrices,
parameterised by the scalar $\kappa_t(r)$.
\end{theorem}

\begin{proof}
Equation~\eqref{eq:H-ray} is just a rearrangement of~\eqref{eq:H-explicit}.
For $M_\rho$: $\det M_\rho = \rho^2 - \rho^2 = 0$, so rank $\le 1$, and the
trace is $1 + \rho^2 > 0$, so rank is exactly $1$. The non-zero eigenvalue
is the trace $1 + \rho^2$; the corresponding eigenvector is $n = (\rho,-1)^\top$
(check: $M_\rho n = (\rho^3+\rho, -\rho^2 -1)^\top = (1+\rho^2)(\rho,-1)^\top$).
$\Box$
\end{proof}

\begin{interpretation}
The Gaussian case has $H_t = \Sigma_t^{-1}$ constant in $x$: the curvature
field is a single matrix. The Laplace case has $H_t$ varying with $x$, but
only along the $1$D ray indexed by the scalar $\kappa_t(r)$. This is the
non-Gaussian generalisation of ``precision matrix'': the field is not a
single matrix, but a smooth curve of matrices with rank-1 update, governed
by the scalar function $r\mapsto \kappa_t(r)$.
\end{interpretation}

\section{Alternative expression for $\kappa_t$: the ODE route}

The formula $\kappa_t(r) = -\psi_t'(r)$ can be turned into an explicit
expression that avoids differentiating the tanh-type formula~\eqref{eq:psi-closed}.
The trick is the identity
\begin{equation}
\kappa_t(r) = -\frac{\mathrm{d}^2}{\mathrm{d}r^2}\log h_t(r)
            = -\frac{h_t''(r)}{h_t(r)} + \left(\frac{h_t'(r)}{h_t(r)}\right)^2
            = \psi_t(r)^2 - \frac{h_t''(r)}{h_t(r)}.
\label{eq:kappa-as-psi-sq-plus}
\end{equation}
We will show in Chapter~\ref{ch:fourier} that $h_t$ satisfies the linear
ODE
\begin{equation}
(1 - \btilde^2 \partial_r^2) h_t(r) = \phi_\tau(r),\qquad
\phi_\tau(r) := \frac{1}{\sqrt{2\pi\tau^2}}\exp\!\left(-\frac{r^2}{2\tau^2}\right).
\label{eq:h-ODE}
\end{equation}
Solving for $h_t''$,
\begin{equation}
h_t''(r) = \frac{h_t(r) - \phi_\tau(r)}{\btilde^2}.
\label{eq:hprimeprime}
\end{equation}
Plugging into~\eqref{eq:kappa-as-psi-sq-plus},
\begin{equation}
\boxed{\;\kappa_t(r) = \psi_t(r)^2 - \frac{1}{\btilde^2} + \frac{1}{\btilde^2}\,\frac{\phi_\tau(r)}{h_t(r)}.\;}
\label{eq:kappa-alt}
\end{equation}
This is the form we will use to prove the Mills-ratio closed form of
$\kappa_t(0)$ in Chapter~\ref{ch:mills}. It is also the form most
convenient for numerical evaluation, since the tanh-type formula for
$\psi_t$ is stable and~\eqref{eq:kappa-alt} avoids differentiating it.

\begin{remark}[Sanity check at the tails]
As $r\to\pm\infty$: $\psi_t(r)\to\mp 1/\btilde$, so $\psi_t^2\to 1/\btilde^2$,
cancelling the $-1/\btilde^2$. Meanwhile $\phi_\tau(r)/h_t(r)\to 0$ because
$h_t$ has exponential Laplace tails whereas $\phi_\tau$ has Gaussian tails
that decay strictly faster. Therefore $\kappa_t(\pm\infty) = 0$, consistent
with the saturation of $\psi_t$.
\end{remark}


% ============================================================
\part{Fourier representation}
% ============================================================

The results of Part II were obtained in real space: convolution in
Section~\ref{ch:convolution}, chain rule in Chapter~\ref{ch:score},
linear algebra in Chapter~\ref{ch:curv}. The Fourier picture tells the
same story but in a cleaner geometric language: the characteristic
function of $p_t$ factorises explicitly, the non-Gaussian content
collapses onto a single innovation-frequency axis, and the residual
density satisfies a genuinely linear ODE. We now derive all of this
from scratch.

% ============================================================
\chapter{Characteristic function of the $K{=}1$ law}
\label{ch:fourier}
% ============================================================

\section{Clean characteristic function}

The clean joint law of $(A_0, A_1)$ has characteristic function
\begin{equation}
\hatfn{p}_0(q_0, q_1) := \E\!\left[ e^{i(q_0 A_0 + q_1 A_1)} \right]
 = \E\!\left[ e^{i q_0 A_0 + i q_1 (\alpha A_0 + \varepsilon)} \right]
 = \E\!\left[ e^{i(q_0 + \alpha q_1) A_0} \right]\,
   \E\!\left[ e^{i q_1 \varepsilon} \right],
\label{eq:clean-cf}
\end{equation}
where we used independence of $A_0$ and $\varepsilon$. The two
expectations are standard:
\begin{align}
\E\!\left[ e^{i(q_0+\alpha q_1) A_0} \right]
 &= \exp\!\left( i\mu_0 (q_0+\alpha q_1) - \tfrac{\sigma_0^2}{2}(q_0+\alpha q_1)^2\right), \\
\E\!\left[ e^{i q_1 \varepsilon} \right]
 &= \frac{1}{1 + b^2 q_1^2}, \qquad \varepsilon \sim \Lap(0,b).
\end{align}

\begin{whybox}
The characteristic function of $\Lap(0,b)$ is
\[
\int_\R \frac{1}{2b} e^{-|u|/b}\,e^{iqu}\,\mathrm{d}u
 = \frac{1}{2b}\left[ \int_0^\infty e^{(iq-1/b)u}\,\mathrm{d}u
                     + \int_{-\infty}^0 e^{(iq+1/b)u}\,\mathrm{d}u \right]
 = \frac{1}{1 + b^2 q^2}.
\]
The rational form $1/(1+b^2 q^2)$ is what makes Laplace such a
tractable non-Gaussian noise: it is the simplest nontrivial example
of an ``exponential of a rational pole'' structure in Fourier space.
\end{whybox}

Collecting,
\begin{equation}
\hatfn{p}_0(q_0, q_1) = \frac{1}{1 + b^2 q_1^2}\,
   \exp\!\left(i\mu_0 (q_0+\alpha q_1) - \tfrac{\sigma_0^2}{2}(q_0+\alpha q_1)^2\right).
\label{eq:clean-cf-final}
\end{equation}

\section{OU corruption in Fourier space}

The channel $X_k = \beta A_k + \sqrt{\Delta} Z_k$ with
$Z_k \perp A$ implies, for the joint characteristic function
$\hatfn{p}_t(k_0, k_1) := \E[e^{i(k_0 X_0 + k_1 X_1)}]$:
\begin{align}
\hatfn{p}_t(k_0,k_1)
 &= \E\!\left[ e^{i(k_0 (\beta A_0 + \sqrt{\Delta} Z_0) + k_1(\beta A_1 + \sqrt{\Delta}Z_1))}\right] \nonumber\\
 &= \E\!\left[ e^{i\beta(k_0 A_0 + k_1 A_1)}\right]\cdot
    \E\!\left[ e^{i\sqrt{\Delta}(k_0 Z_0 + k_1 Z_1)}\right] \nonumber\\
 &= \hatfn{p}_0(\beta k_0, \beta k_1)\,\cdot\, e^{-\Delta(k_0^2 + k_1^2)/2},
\label{eq:cf-OU-corrupt}
\end{align}
using independence of $A$ from $(Z_0, Z_1)$ and the standard Gaussian
CF $\E[e^{iqZ}] = e^{-q^2/2}$. Substituting~\eqref{eq:clean-cf-final}
with $(q_0,q_1) = (\beta k_0, \beta k_1)$ and writing
$\btilde = \beta b$,
\begin{align}
\hatfn{p}_t(k_0, k_1)
 &= \frac{1}{1+\btilde^2 k_1^2}\,
    \exp\!\left(i\beta\mu_0(k_0+\alpha k_1)
               - \tfrac{\beta^2\sigma_0^2}{2}(k_0+\alpha k_1)^2
               - \tfrac{\Delta}{2}(k_0^2+k_1^2)\right).
\label{eq:cf-Kone}
\end{align}

\begin{theorem}[Fourier representation, $K{=}1$]
\label{thm:fourier-K1}
The characteristic function of the noisy joint law is given
by~\eqref{eq:cf-Kone}. In the rotated innovation-frequency coordinates
\begin{equation}
\ell_0 := k_0 + \alpha k_1, \qquad
\ell_1 := k_1,
\label{eq:ell-rotation-K1}
\end{equation}
(equivalently, $k_0 = \ell_0 - \alpha \ell_1$, $k_1 = \ell_1$),
the characteristic function has the form
\begin{equation}
\hatfn{q}_t(\ell_0, \ell_1) :=
\hatfn{p}_t(\ell_0 - \alpha\ell_1, \ell_1)
= \underbrace{e^{i\beta\mu_0 \ell_0
                - \frac{\beta^2\sigma_0^2}{2}\ell_0^2
                - \frac{\Delta}{2}[(\ell_0 - \alpha\ell_1)^2 + \ell_1^2]}}_{\text{Gaussian sector}}
  \cdot \underbrace{\frac{1}{1 + \btilde^2 \ell_1^2}}_{\text{rational sector}}.
\label{eq:cf-rotated-K1}
\end{equation}
The non-Gaussian content is a function of $\ell_1$ alone; the
Gaussian sector is a nondegenerate quadratic form in $(\ell_0,\ell_1)$.
\end{theorem}

\begin{proof}
Equation~\eqref{eq:cf-Kone} is derived above.
Substitute $k_0 + \alpha k_1 = \ell_0$ and $k_1 = \ell_1$;
also $k_0^2 + k_1^2 = (\ell_0 - \alpha\ell_1)^2 + \ell_1^2$. The rational
factor becomes $1/(1 + \btilde^2 \ell_1^2)$. No further simplification is
needed.
\end{proof}

\begin{interpretation}[Why this rotation is natural]
The clean chain has the causal decomposition
$A_0, \varepsilon = A_1 - \alpha A_0$ with $A_0, \varepsilon$ independent.
Their Fourier duals are $k_0 + \alpha k_1$ (the conjugate of $A_0$) and
$k_1$ (the conjugate of $\varepsilon$). The rotation~\eqref{eq:ell-rotation-K1}
is \emph{exactly} this causal coordinate change, rewritten in Fourier space.
After the rotation, the Laplace innovation contributes one rational factor
in $\ell_1$ only, decoupled from the site frequency $\ell_0$.
\end{interpretation}

\section{From characteristic function to residual ODE}
\label{sec:residual-ODE}

We now derive~\eqref{eq:h-ODE}, the linear ODE satisfied by the residual
density $h_t(r)$. Start from the observation that $R_t = \beta\varepsilon + G_t$
has characteristic function
\begin{equation}
\hatfn{h}_t(\ell) := \E[e^{i\ell R_t}]
 = \E[e^{i\ell\beta\varepsilon}]\,\E[e^{i\ell G_t}]
 = \frac{1}{1 + \btilde^2 \ell^2}\cdot e^{-\tau^2 \ell^2/2}.
\label{eq:htilde-lit}
\end{equation}
Multiply both sides by $(1 + \btilde^2 \ell^2)$:
\begin{equation}
(1 + \btilde^2 \ell^2)\,\hatfn{h}_t(\ell) = e^{-\tau^2 \ell^2/2}.
\label{eq:htilde-clean}
\end{equation}
Invert the Fourier transform. Under the convention
$h_t(r) = \frac{1}{2\pi}\int e^{-i\ell r}\hatfn{h}_t(\ell)\,\mathrm{d}\ell$,
multiplication by $i\ell$ corresponds to $\partial_r$; in particular
$(i\ell)^2 = -\ell^2$ corresponds to $\partial_r^2$, so multiplication by
$\ell^2$ corresponds to $-\partial_r^2$. Applying the inverse transform to
both sides of~\eqref{eq:htilde-clean}:
\begin{equation}
(1 - \btilde^2 \partial_r^2)\,h_t(r) = \phi_\tau(r),\qquad
\phi_\tau(r) := \frac{1}{\sqrt{2\pi\tau^2}} e^{-r^2/(2\tau^2)}.
\label{eq:h-ODE-repeat}
\end{equation}
This is~\eqref{eq:h-ODE}, rigorously derived.

\begin{remark}[Check: the clean limit]
At $t=0$: $\tau = 0$, so $\phi_\tau(r)$ becomes a Dirac delta $\delta(r)$;
the ODE reduces to $(1 - b^2 \partial_r^2) h_0 = \delta$, whose solution
(with decay at infinity) is the Green's function
$h_0(r) = (1/2b)e^{-|r|/b}$, the clean Laplace density. Good.
\end{remark}

\begin{remark}[Check: the far-$t$ limit]
As $t\to\infty$: $\btilde\to 0$, so the ODE reduces to
$h_\infty(r) = \phi_\tau(r)$ (a pure Gaussian with variance $\tau^2\to 1$).
This recovers the Gaussian benchmark in the fully-noised regime.
\end{remark}

\section{Recovering $h_t$ from the ODE: an alternative route}

The ODE~\eqref{eq:h-ODE-repeat} gives an alternative way of arriving at
the closed form~\eqref{eq:h-closed}, bypassing the direct convolution.
The Green's function of $1 - \btilde^2 \partial_r^2$ that decays at
infinity is $(2\btilde)^{-1} e^{-|r|/\btilde}$, so
\begin{equation}
h_t(r) = \int_\R \frac{1}{2\btilde}\,e^{-|r-u|/\btilde}\,\phi_\tau(u)\,\mathrm{d}u.
\end{equation}
This is not our original convolution, but a dual one: a Laplace density
convolved against a Gaussian. By the symmetry of convolution, however,
this is the same object, and working through the completion-of-squares
in $u$ (instead of in $r$) gives~\eqref{eq:h-closed} with the roles of
$L_+$ and $L_-$ symmetrised.

\begin{whybox}
Why the ODE matters: from~\eqref{eq:h-ODE-repeat} we read off
$h_t'' = (h_t - \phi_\tau)/\btilde^2$ without having to differentiate
\eqref{eq:h-closed} twice. This is the single fact that made the
alternative expression~\eqref{eq:kappa-alt} for $\kappa_t(r)$ possible --
and that alternative expression is exactly what delivers the Mills-ratio
closed form at $r=0$ in the next part.
\end{whybox}

\section{A diagnostic: where nonlinearity enters}
\label{sec:nonlinearity-diagnosis}

In Fourier space, everything is \emph{linear}: the characteristic
function, the rational rational factor, the ODE. The \emph{nonlinearity}
of the score and curvature comes from a single operation: taking the
logarithm of the real-space density.
\begin{itemize}[leftmargin=2em, topsep=3pt]
\item $h_t$ is a linear combination of exponential-times-Gaussian-CDF
terms.
\item $\log h_t$ is not linear in those terms; it combines them via a
log-sum-exp, which gives a tanh-type saturating score.
\item $\partial_r \log h_t = h_t'/h_t$ is therefore bounded (saturates);
$\partial_r^2 \log h_t$ is concentrated near $r=0$ (where the ratio
$\phi_\tau/h_t$ peaks) and vanishes at infinity.
\end{itemize}
This is the cleanest identification of where non-Gaussianity lives in the
problem. It foreshadows the broader story of the thesis: \emph{the density
is linearly structured in Fourier; the score is where nonlinearity
takes hold, and it is concentrated in the innovation direction.}


% ============================================================
\part{Scalar observables}
% ============================================================

The structure theorems of Part II describe the entire residual density
$h_t(r)$, the entire score $\psi_t(r)$, and the entire curvature field
$\kappa_t(r)$. These are \emph{functions} of $r$; to get finite-dimensional
diagnostics that we can track as $t$ varies -- and that play the role of
the Gaussian precision-matrix mass -- we extract scalar summaries. Each
one is a functional of $h_t$. Our goal is to compute each in closed form
and to interpret what it measures.

% ============================================================
\chapter{Centre curvature: the Mills-ratio closed form}
\label{ch:mills}
% ============================================================

The most important scalar summary is the curvature at the origin,
$\kappa_t(0)$. We state and prove in full the identity that expresses it
through Mills' ratio, and then interpret its limits.

\section{Statement}

\begin{theorem}[Closed form for $\kappa_t(0)$]
\label{thm:kappa-center}
Let $\phi, \Phi$ denote the standard normal pdf and cdf, and let
\begin{equation}
\mills(a) := \frac{\Phi(-a)}{\phi(a)}, \qquad a\in\R,
\label{eq:mills-def}
\end{equation}
be Mills' ratio of the standard normal. Then, with $a = \tau/\btilde$,
\begin{equation}
\boxed{\; \kappa_t(0)\,\btilde^2 = \frac{1}{a\,\mills(a)} - 1. \;}
\label{eq:mills-main}
\end{equation}
Equivalently,
\begin{equation}
\kappa_t(0) = \frac{1}{\btilde^2}\left(\frac{1}{a\mills(a)} - 1\right).
\end{equation}
\end{theorem}

\begin{whybox}
Why this is the right object: $\kappa_t(0)\,\btilde^2$ is dimensionless
(both sides have dimension of length squared, canceling). It depends on
$t$ only through the single scalar $a(t) = \tau(t)/\btilde(t)$, which
monotonically increases from $0$ at $t=0$ to $+\infty$ as $t\to\infty$.
Thus the entire one-parameter family of residual densities is compressed
into a single dimensionless curve.
\end{whybox}

\section{Proof}

The proof has three elementary ingredients: the alternative formula for
$\kappa_t$ from Section~\ref{ch:curv}, the closed form for $h_t(0)$, and
Mills' ratio at a single point.

\paragraph{Step 1: $\psi_t(0) = 0$.}
By direct inspection of~\eqref{eq:h-closed}, $h_t(-r) = h_t(r)$: swap the
two terms using $\Phi(r/\tau - a) \leftrightarrow \Phi(-r/\tau - a)$ under
$r\mapsto -r$. So $h_t$ is an even function, and its derivative at $0$ is
zero:
\begin{equation}
\psi_t(0) = \frac{h_t'(0)}{h_t(0)} = 0.
\end{equation}

\paragraph{Step 2: the alternative formula evaluated at $0$.}
From~\eqref{eq:kappa-alt},
\begin{equation}
\kappa_t(0) = \psi_t(0)^2 - \frac{1}{\btilde^2} + \frac{1}{\btilde^2}\,\frac{\phi_\tau(0)}{h_t(0)}
            = -\frac{1}{\btilde^2} + \frac{1}{\btilde^2}\,\frac{\phi_\tau(0)}{h_t(0)}.
\label{eq:mills-step2}
\end{equation}
It remains to evaluate the ratio $\phi_\tau(0)/h_t(0)$.

\paragraph{Step 3: compute $h_t(0)$.}
From~\eqref{eq:h-closed} at $r = 0$:
\begin{align}
h_t(0) &= \frac{e^{a^2/2}}{2\btilde}\left[L_+(0) + L_-(0)\right]
       = \frac{e^{a^2/2}}{2\btilde}\left[ e^0\Phi(-a) + e^0 \Phi(-a)\right]
       = \frac{e^{a^2/2}}{\btilde}\Phi(-a).
\label{eq:h0-value}
\end{align}
(We used $L_+(0) = \Phi(-a)$ and $L_-(0) = \Phi(-a)$ from~\eqref{eq:Lpm-def}.)

\paragraph{Step 4: compute $\phi_\tau(0)/h_t(0)$.}
$\phi_\tau(0) = 1/\sqrt{2\pi\tau^2}$, and rewriting using $\tau = a\btilde$:
\begin{equation}
\phi_\tau(0) = \frac{1}{\sqrt{2\pi}\, a\btilde}.
\end{equation}
Dividing,
\begin{equation}
\frac{\phi_\tau(0)}{h_t(0)}
 = \frac{1/(\sqrt{2\pi}\, a\btilde)}{e^{a^2/2}\Phi(-a)/\btilde}
 = \frac{1}{a\,\sqrt{2\pi}\,e^{a^2/2}\,\Phi(-a)}
 = \frac{\phi(a)}{a\,\phi(a)\,\sqrt{2\pi}\,e^{a^2/2}\,\Phi(-a)}.
\end{equation}
But $\phi(a) = (2\pi)^{-1/2} e^{-a^2/2}$, so
$(2\pi)^{1/2}e^{a^2/2} = 1/\phi(a)$, and
\begin{equation}
\frac{\phi_\tau(0)}{h_t(0)}
 = \frac{\phi(a)}{a\,\Phi(-a)}
 = \frac{1}{a\,\mills(a)},
\label{eq:ratio-is-mills}
\end{equation}
where we used~\eqref{eq:mills-def}.

\paragraph{Step 5: assemble.}
Substituting~\eqref{eq:ratio-is-mills} into~\eqref{eq:mills-step2},
\begin{equation}
\kappa_t(0)\,\btilde^2 = -1 + \frac{1}{a\,\mills(a)}
                       = \frac{1}{a\,\mills(a)} - 1,
\end{equation}
which is~\eqref{eq:mills-main}. $\Box$

\section{Sanity check: the two extreme limits}

Mills' ratio has well-known expansions at $a\to 0$ and $a\to\infty$.
Using them, we recover the clean Laplace and Gaussian limits.

\paragraph{Clean limit $t\to 0^+$ ($a\to 0^+$).}
At $a = 0$: $\phi(0) = 1/\sqrt{2\pi}$, $\Phi(0) = 1/2$, so
$\mills(0) = \sqrt{2\pi}/2 = \sqrt{\pi/2}$.
Expanding,
\begin{equation}
\mills(a) = \sqrt{\pi/2}\,\bigl(1 - a\sqrt{2/\pi} + O(a^2)\bigr),\qquad a\to 0.
\label{eq:mills-small}
\end{equation}
Hence
\begin{equation}
\frac{1}{a\mills(a)} = \frac{1}{a\sqrt{\pi/2}}\bigl(1 + a\sqrt{2/\pi} + O(a^2)\bigr)
 = \frac{\sqrt{2/\pi}}{a} + 1 + O(a).
\end{equation}
Therefore
\begin{equation}
\kappa_t(0)\,\btilde^2 \sim \frac{\sqrt{2/\pi}}{a} \to +\infty \quad (a\to 0^+),
\label{eq:kappa-small-t}
\end{equation}
i.e.\ the centre curvature diverges as $1/a = \btilde/\tau$ when $t\to 0$.
This is the $r=0$ Laplace cusp: the clean Laplace density has a corner
at $0$, so the log-density has infinite negative curvature there.

\paragraph{Fully-noised limit $t\to\infty$ ($a\to\infty$).}
For large $a$,
\begin{equation}
\mills(a) \sim \frac{1}{a} - \frac{1}{a^3} + O(a^{-5}),
\label{eq:mills-large}
\end{equation}
hence $a\mills(a) \sim 1 - 1/a^2 + \cdots$, and
$1/(a\mills(a)) \sim 1 + 1/a^2$. Subtracting 1:
\begin{equation}
\kappa_t(0)\,\btilde^2 \sim \frac{1}{a^2} = \frac{\btilde^2}{\tau^2}.
\end{equation}
Thus $\kappa_t(0) \sim 1/\tau^2$. Compare to the Gaussian-benchmark
curvature $\kappa_t^{(G)} = 1/\tau_G^2 = 1/(\tau^2 + 2\btilde^2)$: for
large $a$, $\tau^2 \gg \btilde^2$, and both agree to leading order.

\begin{interpretation}[What the Mills-ratio form buys us]
Equation~\eqref{eq:mills-main} is not just a compact formula: it is the
explicit statement that the centre curvature is a universal function of
a single dimensionless ratio $a$. As $a$ sweeps from $0$ to $\infty$
along diffusion time, $\kappa_t(0)\btilde^2$ sweeps from $+\infty$ to $0$,
passing through every positive value exactly once (monotonicity is not
proved here but follows from the fact that $a\mills(a)$ is strictly
increasing). Every finite-time non-Gaussianity of the centre of the
density is encoded in this single scalar curve, and the Gaussian
benchmark simply reads as a different universal function of $a$,
differing by a correction of order $1/a^4$ at large $a$.
\end{interpretation}

% ============================================================
\chapter{Gaussian benchmark and the centre-benchmark gap}
\label{ch:benchmark}
% ============================================================

To calibrate non-Gaussianity, we compare the Laplace benchmark against a
\emph{matched-variance} Gaussian model. This gives a rigorous notion of
``how far is the true curvature from the Gaussian fit of the same
two-point statistics.''

\section{The matched Gaussian innovation}

Replace $\varepsilon \sim \Lap(0,b)$ by $\varepsilon^{(G)} \sim \Normal(0, 2b^2)$.
Since $\V(\varepsilon) = 2b^2 = \V(\varepsilon^{(G)})$, the two innovations
have identical second-order statistics. In the Laplace benchmark,
$R_t = \beta\varepsilon + G_t$ has variance $2\btilde^2 + \tau^2 =: \tau_G^2$;
in the matched Gaussian benchmark, the corresponding residual
$R_t^{(G)}$ has the same variance but is Gaussian:
\begin{equation}
R_t^{(G)} \sim \Normal(0, \tau_G^2),\qquad \tau_G^2 := 2\btilde^2 + \tau^2.
\label{eq:tauG2}
\end{equation}
The joint noisy law is jointly Gaussian and the curvature is constant
in $x$:
\begin{equation}
H_t^{(G)} = \begin{pmatrix}1/v + \rho^2/\tau_G^2 & -\rho/\tau_G^2 \\
                         -\rho/\tau_G^2 & 1/\tau_G^2\end{pmatrix}
= \begin{pmatrix}1/v & 0\\ 0 & 0\end{pmatrix} + \kappa_t^{(G)} M_\rho,
\qquad
\kappa_t^{(G)} := \frac{1}{\tau_G^2}.
\label{eq:kappaG}
\end{equation}

\begin{whybox}
Matching $\V(\varepsilon) = \V(\varepsilon^{(G)}) = 2b^2$ means the two
benchmarks cannot be told apart by any two-point statistic of the
residual. Everything that distinguishes them -- the cusp at $r=0$, the
saturating score, the location-dependent curvature -- is by construction
higher-order.
\end{whybox}

\section{The centre-benchmark gap}

The residual curvature at $r=0$ of the Laplace model is
$\kappa_t(0) = (1/\btilde^2)(1/(a\mills(a)) - 1)$.
The matched Gaussian gives constant $\kappa_t^{(G)} = 1/(\tau^2 + 2\btilde^2)
= 1/(\btilde^2(a^2 + 2))$. Their gap is
\begin{equation}
\delta\kappa_t := \kappa_t(0) - \kappa_t^{(G)}
= \frac{1}{\btilde^2}\left[\frac{1}{a\mills(a)} - 1 - \frac{1}{a^2 + 2}\right].
\label{eq:gap-exact}
\end{equation}

\paragraph{Small $a$ (small $t$): the cusp dominates.}
$1/(a\mills(a)) \sim \sqrt{2/\pi}/a$ and $1/(a^2+2) \to 1/2$, so
$\delta\kappa_t \sim (\btilde^2)^{-1}\sqrt{2/\pi}/a$, i.e.\ diverges
like $1/a$. The Gaussian fit misses the Laplace cusp catastrophically
in the clean regime.

\paragraph{Large $a$ (large $t$): the gap closes.}
Use $1/(a\mills(a)) \sim 1 + 1/a^2 + 3/a^4 + \cdots$ (from the expansion
$\mills(a) = 1/a - 1/a^3 + 3/a^5 - \cdots$) and
$1/(a^2 + 2) = 1/a^2 - 2/a^4 + O(a^{-6})$:
\begin{align}
\frac{1}{a\mills(a)} - 1 - \frac{1}{a^2+2}
 &= \left(\frac{1}{a^2} + \frac{3}{a^4} + O(a^{-6})\right)
    - \left(\frac{1}{a^2} - \frac{2}{a^4} + O(a^{-6})\right) \nonumber \\
 &= \frac{5}{a^4} + O(a^{-6}).
\end{align}
So $\delta\kappa_t \sim 5/(\btilde^2 a^4) = 5\btilde^2/\tau^4$, which
goes to zero as $t\to\infty$. The gap closes quartically.

\begin{interpretation}[Where the Gaussian benchmark is honest]
For small $t$ (strong non-Gaussianity), the Gaussian benchmark
\emph{underestimates} centre curvature by a factor growing like $\btilde/\tau$;
a score trained against a Gaussian fit will miss the cusp entirely.
For large $t$, the gap closes extremely fast ($\btilde^2/\tau^4$), so
non-Gaussian fine-tuning yields negligible benefit in the
late-diffusion regime. This is the quantitative underpinning of the
folklore that diffusion models ``don't need to care about non-Gaussianity
at late times.''
\end{interpretation}

% ============================================================
\chapter{Kurtosis and bandwidth}
\label{ch:kurtosis}
% ============================================================

\section{The fourth cumulant of $R_t$}

\begin{proposition}[Residual fourth cumulant]
\label{prop:kurt}
Let $\kappa_4(Y)$ denote the fourth cumulant of a real random variable
$Y$. Then
\begin{equation}
\kappa_4(R_t) = 12\,\beta^4\,b^4,
\label{eq:kappa4}
\end{equation}
and consequently the excess kurtosis of $R_t$ is
\begin{equation}
\frac{\kappa_4(R_t)}{\V(R_t)^2} = \frac{12\beta^4 b^4}{(2\btilde^2 + \tau^2)^2}
 = \frac{12 \btilde^4}{\btilde^4(a^2+2)^2}
 = \frac{12}{(a^2+2)^2}
 = \frac{3}{(1 + a^2/2)^2}.
\label{eq:excess-kurt}
\end{equation}
\end{proposition}

\begin{proof}
$R_t = \beta\varepsilon + G_t$ is a sum of independent variables, so
cumulants add: $\kappa_4(R_t) = \kappa_4(\beta\varepsilon) + \kappa_4(G_t)$.
The Gaussian $G_t$ has $\kappa_4 = 0$.
The fourth cumulant of $\beta\varepsilon$ scales as $\beta^4$ times that of
$\varepsilon \sim \Lap(0,b)$. For Laplace$(0,b)$, the fourth moment is
$24 b^4$ and the second moment is $2b^2$, so
$\kappa_4 = m_4 - 3 m_2^2 = 24b^4 - 3(2b^2)^2 = 12b^4$. (Here
$\kappa_3 = 0$ by symmetry.) Hence
$\kappa_4(R_t) = 12\beta^4 b^4$. Substituting $\V(R_t) = 2\btilde^2 + \tau^2$
and $\btilde = \beta b$ gives~\eqref{eq:excess-kurt}. $\Box$
\end{proof}

\begin{interpretation}[Non-Gaussianity persistence]
At $t=0$ ($a=0$): excess kurtosis is $3$, the Laplace value.
At $t\to\infty$ ($a\to\infty$): excess kurtosis $\to 0$, the Gaussian
limit. The decay is $\sim 12/a^4$, the same quartic rate as the
centre-benchmark gap. This is not a coincidence: both quantities are
measuring the same thing -- the fourth-order departure from Gaussianity
of $R_t$ -- just through different functionals (one pointwise at $r=0$,
one integrated).
\end{interpretation}

\section{Full width at half maximum of $\kappa_t$}

Since $\kappa_t(r)$ is an even function peaked at $r=0$ and decaying to
zero at $\pm\infty$, the characteristic width of the curvature
band is
\begin{equation}
\mathrm{FWHM}(t) := 2 r^*(t), \qquad \kappa_t(r^*) = \tfrac{1}{2}\kappa_t(0).
\end{equation}
This is the non-Gaussian generalisation of the ``tridiagonal mass''
bandwidth $b_{0.95}(t)$ used in the Gaussian case: it quantifies how
localised the curvature is to the residual-ray centre.

\begin{whybox}
Why FWHM rather than another width: the function $\kappa_t(r)$ has
heavy tails (they decay to zero) but a clear peak at $0$. Second moment
is dominated by the tails and does not measure concentration; variance
of the corresponding curvature mass is ill-defined; FWHM is
well-defined, robust, and easy to compute from the closed form.
\end{whybox}

\paragraph{Two natural length scales.}
The residual density $h_t(r)$ is built from two elementary scales:
$\btilde$ (Laplace scale of the rescaled innovation) and $\tau$
(Gaussian scale of the smoothing). Numerical experiments (Figure 4 of
the main memo) confirm:
\begin{itemize}[leftmargin=2em, topsep=3pt]
\item $\mathrm{FWHM}(t) \approx \btilde$ for small $t$ (Laplace regime, $a\ll 1$);
\item $\mathrm{FWHM}(t) \approx \tau$ for large $t$ (Gaussian regime, $a\gg 1$);
\item crossover around $a = \tau/\btilde \approx 1$.
\end{itemize}
A closed form for $\mathrm{FWHM}(t)$ is possible in principle via the
alternative formula~\eqref{eq:kappa-alt} but is not illuminating;
the scales $\btilde$ and $\tau$ provide the right intuition.


% ============================================================
\part{The $K{=}2$ case}
% ============================================================

The $K=1$ case is a pedagogical exercise: one innovation, one residual,
one non-Gaussian scalar. The $K=2$ case is where the central question of
the thesis starts to bite. Two innovations are coupled through the
shared latent state $A_1$; can we still factorise the exact density into
a product of one-dimensional non-Gaussian bonds? The answer, as we will
show in full, is \emph{no}: there is a specific Gaussian
cross-covariance, $-\alpha\Delta$, that forbids this factorisation at
any finite $t$. But the \emph{Markov locality pattern} survives:
non-Gaussianity remains diagonal in innovation coordinates, and the
Gaussian sector only couples nearest neighbours. We work out every step.

% ============================================================
\chapter{Characteristic function for $K{=}2$}
\label{ch:K2-fourier}
% ============================================================

\section*{A note on rotation conventions}

Before diving in, we flag a convention change relative to Part III.
Two different rotations of the frequency variables are natural in this
problem, and we use each where it is cleanest.

\paragraph{Causal (prior/innovation) rotation.} This is the Fourier dual
of the real-space innovation coordinates
$(x_0, y_1, y_2, \ldots)$ with $y_k = x_k - \alpha x_{k-1}$.
It diagonalises the linear phase of the characteristic function
(making it depend only on the first coordinate), but it leaves the
Laplace rational factors on non-axis-aligned directions.

\paragraph{Minimal (innovation) rotation.}
This is the rotation that aligns the Laplace rational factors with the
coordinate axes, but leaves the linear phase in a general direction.

\medskip

At $K=1$ these two rotations coincide (there is only one innovation),
and Part III uses the causal convention
$\ell_0 = k_0 + \alpha k_1,\ \ell_1 = k_1$. This choice highlights the
one-dimensional non-Gaussian axis $\ell_1$.

At $K=2$ the two rotations differ. We will use the \emph{minimal}
convention $\ell_0 = k_0,\ \ell_1 = k_1 + \alpha k_2,\ \ell_2 = k_2$
because it is the one that separates the two Laplace innovations onto
distinct axes --- and separating the innovations is precisely the
content we need to test when asking whether the $K=1$ factorisation
extends. A reader who prefers the fully-causal rotation will find it
equivalent (the two differ by a further Gaussian shear in the site axis),
and will obtain a slightly different $M$ matrix whose off-diagonal
structure is tridiagonal rather than having one block-column pattern;
the cross-term $M_{12} = -\alpha\Delta$ is the same in both conventions
and is what matters for the obstruction theorem.

\section{Clean characteristic function}

The clean chain for $K = 2$ has
\begin{equation}
A_0 \sim \Normal(\mu_0, \sigma_0^2),\qquad
A_1 = \alpha A_0 + \varepsilon_1,\qquad
A_2 = \alpha A_1 + \varepsilon_2,
\end{equation}
with $\varepsilon_1, \varepsilon_2 \stackrel{\mathrm{iid}}{\sim} \Lap(0, b)$
independent of $A_0$.

The clean characteristic function is
\begin{equation}
\hatfn p_0(q_0, q_1, q_2) =
\E\!\left[ e^{i(q_0 A_0 + q_1 A_1 + q_2 A_2)}\right].
\end{equation}
Substitute $A_1 = \alpha A_0 + \varepsilon_1$ and
$A_2 = \alpha^2 A_0 + \alpha\varepsilon_1 + \varepsilon_2$:
\begin{equation}
q_0 A_0 + q_1 A_1 + q_2 A_2
 = (q_0 + \alpha q_1 + \alpha^2 q_2) A_0 + (q_1 + \alpha q_2)\varepsilon_1
 + q_2\varepsilon_2.
\end{equation}
By independence of $A_0, \varepsilon_1, \varepsilon_2$, the characteristic
function factorises:
\begin{equation}
\hatfn p_0(q_0, q_1, q_2) =
\underbrace{\E[e^{iq A_0}]}_{q = q_0+\alpha q_1+\alpha^2 q_2}\,
\underbrace{\E[e^{i(q_1+\alpha q_2)\varepsilon_1}]}_{=1/(1+b^2(q_1+\alpha q_2)^2)}\,
\underbrace{\E[e^{iq_2\varepsilon_2}]}_{=1/(1+b^2 q_2^2)}.
\label{eq:K2-clean-cf-pre}
\end{equation}
With the Gaussian CF $\E[e^{iq A_0}] = \exp(i\mu_0 q - \sigma_0^2 q^2/2)$,
\begin{equation}
\boxed{\;
\hatfn p_0(q_0, q_1, q_2)
 = \frac{e^{i\mu_0 q - \sigma_0^2 q^2/2}}
        {(1 + b^2(q_1+\alpha q_2)^2)(1 + b^2 q_2^2)},
 \qquad q := q_0 + \alpha q_1 + \alpha^2 q_2.
\;}
\label{eq:K2-clean-cf}
\end{equation}

\begin{whybox}
Notice that the rational factor has two separate denominators: one for
each innovation. The first, $1 + b^2(q_1 + \alpha q_2)^2$, comes from
$\varepsilon_1 = A_1 - \alpha A_0$; the second, $1 + b^2 q_2^2$, from
$\varepsilon_2 = A_2 - \alpha A_1$. The \emph{rotation} that separates
them is not $(q_0, q_1, q_2) \to (q_0, q_1, q_2)$ but a linear change of
coordinates we will identify below.
\end{whybox}

\section{OU corruption}

The argument of Section~\ref{sec:residual-ODE} extends verbatim:
the joint CF of the noisy chain is
\begin{equation}
\hatfn p_t(k_0, k_1, k_2) = \hatfn p_0(\beta k_0, \beta k_1, \beta k_2)\,
    e^{-\Delta(k_0^2 + k_1^2 + k_2^2)/2}.
\end{equation}
Substituting~\eqref{eq:K2-clean-cf} and writing $\btilde = \beta b$,
$\tilde q := \beta q = \beta(k_0 + \alpha k_1 + \alpha^2 k_2)$:
\begin{align}
\hatfn p_t(k_0, k_1, k_2) =
\frac{\exp\!\left[i\beta\mu_0(k_0+\alpha k_1+\alpha^2 k_2)
       - \tfrac{\beta^2\sigma_0^2}{2}(k_0+\alpha k_1+\alpha^2 k_2)^2
       - \tfrac{\Delta}{2}(k_0^2+k_1^2+k_2^2)\right]}
     {(1 + \btilde^2(k_1+\alpha k_2)^2)(1 + \btilde^2 k_2^2)}.
\label{eq:K2-fourier-original}
\end{align}

\section{Rotation to innovation-frequency coordinates}

Define
\begin{equation}
\ell_0 := k_0,\qquad
\ell_1 := k_1 + \alpha k_2,\qquad
\ell_2 := k_2.
\label{eq:ell-rotation-K2}
\end{equation}
This change of variables is invertible with
\begin{equation}
k_0 = \ell_0,\qquad
k_1 = \ell_1 - \alpha\ell_2,\qquad
k_2 = \ell_2.
\end{equation}
In the new coordinates:

\begin{itemize}[leftmargin=2em, topsep=3pt]
\item The innovation-2 rational factor is $1/(1 + \btilde^2 \ell_2^2)$
(unchanged, since $k_2 = \ell_2$).
\item The innovation-1 rational factor becomes
$1/(1 + \btilde^2 \ell_1^2)$ (exactly, using
$k_1 + \alpha k_2 = \ell_1 - \alpha\ell_2 + \alpha\ell_2 = \ell_1$).
\item The linear phase term $k_0 + \alpha k_1 + \alpha^2 k_2$ becomes
$\ell_0 + \alpha(\ell_1 - \alpha\ell_2) + \alpha^2\ell_2 = \ell_0 + \alpha\ell_1$.
\emph{The $\ell_2$ contribution cancels.}
\end{itemize}

The rational part is now \emph{fully diagonal}: each Laplace factor
depends on one innovation frequency only. The Gaussian sector is a
quadratic form in $(\ell_0, \ell_1, \ell_2)$, which we compute carefully.

\paragraph{Computing the Gaussian quadratic form.}
Let $\mathcal{Q}$ denote the quadratic form appearing in the exponent:
\begin{equation}
\mathcal{Q}(k_0, k_1, k_2) = \beta^2\sigma_0^2 (k_0+\alpha k_1+\alpha^2 k_2)^2
 + \Delta(k_0^2+k_1^2+k_2^2).
\end{equation}
Under the rotation:
\begin{align}
(k_0+\alpha k_1+\alpha^2 k_2)^2 &= (\ell_0 + \alpha\ell_1)^2
  = \ell_0^2 + 2\alpha\ell_0\ell_1 + \alpha^2\ell_1^2, \\
k_0^2 &= \ell_0^2, \\
k_1^2 &= (\ell_1 - \alpha\ell_2)^2 = \ell_1^2 - 2\alpha\ell_1\ell_2 + \alpha^2\ell_2^2, \\
k_2^2 &= \ell_2^2.
\end{align}
Summing,
\begin{align}
\mathcal{Q}(\ell_0, \ell_1, \ell_2) =\;&
\beta^2\sigma_0^2 (\ell_0^2 + 2\alpha\ell_0\ell_1 + \alpha^2\ell_1^2) \nonumber \\
&+ \Delta\left(\ell_0^2 + \ell_1^2 - 2\alpha\ell_1\ell_2 + \alpha^2\ell_2^2 + \ell_2^2\right) \nonumber \\
=\;& \underbrace{(\Delta + \beta^2\sigma_0^2)}_{M_{00}}\ell_0^2
  + \underbrace{(\Delta + \alpha^2\beta^2\sigma_0^2)}_{M_{11}}\ell_1^2
  + \underbrace{\Delta(1+\alpha^2)}_{M_{22}}\ell_2^2 \nonumber \\
& + 2\underbrace{\alpha\beta^2\sigma_0^2}_{M_{01}}\ell_0\ell_1
  + 2\underbrace{(-\alpha\Delta)}_{M_{12}}\ell_1\ell_2 + 2\cdot 0\cdot \ell_0\ell_2.
\label{eq:Q-rotated}
\end{align}
Thus $\mathcal{Q}(\ell) = \ell^\top M \ell$ with
\begin{equation}
\boxed{\;
M = \begin{pmatrix}
\Delta + \beta^2\sigma_0^2 & \alpha\beta^2\sigma_0^2 & 0 \\
\alpha\beta^2\sigma_0^2 & \Delta + \alpha^2\beta^2\sigma_0^2 & -\alpha\Delta \\
0 & -\alpha\Delta & \Delta(1 + \alpha^2)
\end{pmatrix}.
\;}
\label{eq:M-matrix}
\end{equation}

\begin{whybox}
Notice three things in $M$:
(i) $M_{02} = 0$: the site frequency $\ell_0$ does not couple to the
last innovation frequency $\ell_2$. This is the Fourier-space signature of
Markov locality: the chain has only nearest-neighbour coupling.
(ii) $M_{01} = \alpha\beta^2\sigma_0^2 \neq 0$: the site frequency
\emph{does} couple to the first innovation frequency. This is the
prior's effect on the first step of the chain.
(iii) $M_{12} = -\alpha\Delta \neq 0$: the two innovation frequencies
\emph{do} couple. This is the obstruction we now examine.
\end{whybox}

\section{Theorem: the characteristic function factorisation}

\begin{theorem}[Fourier representation, $K{=}2$]
\label{thm:K2-fourier}
In the rotated coordinates~\eqref{eq:ell-rotation-K2},
\begin{equation}
\hatfn q_t(\ell) :=
\hatfn p_t(\ell_0, \ell_1 - \alpha\ell_2, \ell_2)
= e^{i\beta\mu_0(\ell_0 + \alpha\ell_1) - \tfrac{1}{2}\ell^\top M \ell}\cdot
  \frac{1}{1 + \btilde^2 \ell_1^2}\cdot \frac{1}{1 + \btilde^2 \ell_2^2},
\label{eq:K2-rotated-cf}
\end{equation}
with $M$ given by~\eqref{eq:M-matrix}.
\end{theorem}

\begin{proof}
Direct substitution, using the algebra above.
\end{proof}

\begin{remark}[The clean-chain limit]
At $t=0$: $\Delta=0$, $\beta=1$, so
$M = \sigma_0^2\,\mathrm{diag}(1, \alpha^2, 0) + \cdots$ and in fact
$M_{22} = 0$, i.e.\ the Gaussian sector becomes degenerate in the $\ell_2$
direction. This matches the fact that at $t=0$ the chain is
Laplace-innovation-only in $A_2 - \alpha A_1 = \varepsilon_2$, with no
Gaussian smoothing, and the rational factor $1/(1 + b^2 \ell_2^2)$ is
a pure Laplace distribution on its own.
\end{remark}

% ============================================================
\chapter{The obstruction to one-dimensional closure}
\label{ch:K2-obstruction}
% ============================================================

\section{Statement of the obstruction}

For $K=1$, the residual density $h_t(r)$ was a function of a single
coordinate $r = x_1 - \nu - \rho x_0$. Theorem~\ref{thm:K2-fourier}
isolates two innovation frequencies $\ell_1, \ell_2$, whose
\emph{rational} factors are already separated. If we could further
decouple their \emph{Gaussian} contributions, we would get a full
product factorisation into a site mode (in $\ell_0$) and two independent
bonds (in $\ell_1$ and $\ell_2$). The question is whether this is
possible.

\begin{theorem}[$K{=}2$ obstruction]
\label{thm:K2-obstruction}
For any $\alpha \neq 0$ and $t > 0$, the joint noisy law
$p_t(x_0, x_1, x_2)$ does not admit a factorisation of the form
\begin{equation}
p_t(x_0, x_1, x_2) \stackrel{?}{=} g_t(x_0)\, h_t^{(1)}(r_1)\, h_t^{(2)}(r_2),
\label{eq:K2-false-factorisation}
\end{equation}
with $r_1, r_2$ independent residual coordinates and
$h_t^{(1)}, h_t^{(2)}$ one-dimensional densities.
\end{theorem}

\begin{proof}
A factorisation of the form~\eqref{eq:K2-false-factorisation} would, in
Fourier space, correspond to a fully separated characteristic function
\[
\hatfn q_t^{\text{factor}}(\ell_0, \ell_1, \ell_2) =
  \mathcal G_0(\ell_0) \cdot \mathcal B_1(\ell_1) \cdot \mathcal B_2(\ell_2),
\]
where $\mathcal G_0$ is the site factor and $\mathcal B_1, \mathcal B_2$
are the bond factors. This would require in particular the quadratic
form in $(\ell_0, \ell_1, \ell_2)$ to be diagonal after the rotation.
But we have computed $M$ explicitly, and $M_{12} = -\alpha\Delta \neq 0$
whenever $\alpha \neq 0$ and $\Delta \neq 0$ (i.e.\ $t \neq 0$).
Therefore no such factorisation exists. $\Box$
\end{proof}

\begin{interpretation}[The physical origin of $-\alpha\Delta$]
The cross-term $M_{12} = -\alpha\Delta$ comes from the expansion
$k_1^2 = (\ell_1 - \alpha\ell_2)^2 = \ell_1^2 - 2\alpha\ell_1\ell_2 + \alpha^2\ell_2^2$.
It reflects the fact that the OU-injected noise $\sqrt{\Delta}Z_1$ in
coordinate $X_1$ contributes to both the first innovation channel (via
$\ell_1 = k_1 + \alpha k_2$) and the coupling with the second
innovation's dynamics (via the shift $k_1 \to \ell_1 - \alpha\ell_2$).
Markov locality survives (there is no $\ell_0\ell_2$ coupling in $M$),
but the two consecutive innovations pick up a shared dependence on
$Z_1$, the OU noise that lives at the shared vertex $A_1$ of the chain.
\end{interpretation}

\section{Schur complement: effective coupling of the innovations}

Marginalising out $\ell_0$ in the Gaussian sector gives the effective
covariance of $(\ell_1, \ell_2)$ as the Schur complement
\begin{equation}
\widetilde M := M_{(1:2,1:2)} - M_{(1:2,0)}\, M_{00}^{-1}\, M_{(0,1:2)}.
\end{equation}

\paragraph{Computation of the Schur complement.}
Using~\eqref{eq:M-matrix}:
\begin{equation}
M_{(1:2,0)} = \begin{pmatrix} \alpha\beta^2\sigma_0^2 \\ 0 \end{pmatrix},\qquad
M_{00} = \Delta + \beta^2\sigma_0^2 = v,\qquad
M_{(0,1:2)} = \begin{pmatrix} \alpha\beta^2\sigma_0^2 & 0\end{pmatrix}.
\end{equation}
So
\begin{equation}
M_{(1:2,0)}\,M_{00}^{-1}\, M_{(0,1:2)} =
\frac{1}{v}\begin{pmatrix}\alpha^2\beta^4\sigma_0^4 & 0 \\ 0 & 0\end{pmatrix}.
\end{equation}
Therefore
\begin{equation}
\widetilde M = \begin{pmatrix}
\Delta + \alpha^2\beta^2\sigma_0^2 - \frac{\alpha^2\beta^4\sigma_0^4}{v} & -\alpha\Delta \\
-\alpha\Delta & \Delta(1+\alpha^2)
\end{pmatrix}.
\end{equation}
Simplify the $(1,1)$-entry:
\begin{align}
\Delta + \alpha^2\beta^2\sigma_0^2 - \frac{\alpha^2\beta^4\sigma_0^4}{v}
 &= \Delta + \alpha^2\beta^2\sigma_0^2\cdot\frac{v - \beta^2\sigma_0^2}{v}
  = \Delta + \alpha^2\beta^2\sigma_0^2\cdot\frac{\Delta}{v} \nonumber\\
 &= \Delta\left(1 + \frac{\alpha^2\beta^2\sigma_0^2}{v}\right)
  = \frac{\Delta(v + \alpha^2\beta^2\sigma_0^2)}{v}
  = \frac{\Delta(\Delta + \beta^2\sigma_0^2(1+\alpha^2))}{v},
\end{align}
where we used $v = \Delta + \beta^2\sigma_0^2$ twice.

\begin{equation}
\boxed{\;
\widetilde M = \begin{pmatrix}
\dfrac{\Delta(\Delta + \beta^2\sigma_0^2(1+\alpha^2))}{v} & -\alpha\Delta \\[6pt]
-\alpha\Delta & \Delta(1+\alpha^2)
\end{pmatrix},\qquad
\widetilde M_{12} = -\alpha\Delta.
\;}
\label{eq:Schur}
\end{equation}

\begin{remark}
$\widetilde M_{12} = -\alpha\Delta$ is the same cross-term as $M_{12}$. The
marginalisation of $\ell_0$ does not change the $(1,2)$ off-diagonal: the
coupling between $\ell_1$ and $\ell_2$ is \emph{purely} a consequence of
the OU noise at $A_1$, not of any indirect path through the site variable
$A_0$.
\end{remark}

% ============================================================
\chapter{Exact four-term representation}
\label{ch:K2-four-terms}
% ============================================================

We now give an exact real-space representation of $p_t(x_0, x_1, x_2)$ as
a sum of four terms, each involving a bivariate normal CDF. This is the
$K=2$ analogue of the two-term representation~\eqref{eq:h-closed}.

\section{Setup}

By the chain rule, $p_t(x_0,x_1,x_2) = p_t(x_0,x_1)\,p_t(x_2\mid x_0,x_1)$.
By Proposition~\ref{prop:markov-survives}, $X_2 \perp X_0 \mid X_1$, so
\begin{equation}
p_t(x_2\mid x_0,x_1) = p_t(x_2\mid x_1).
\end{equation}
The \emph{second-step} conditional is the analogue of the $K=1$ conditional
but with $A_1$ (a non-Gaussian variable) in the role of the latent state.
This is where the analysis deviates from $K=1$.

\paragraph{Non-Gaussian posterior.}
The analogue of the step ``$A_0\mid X_0$ is Gaussian'' at the second step
is: what is the conditional distribution of $A_1 \mid X_1 = x_1$? Using
Bayes:
\begin{equation}
p_t(A_1 = a \mid X_1 = x_1) \propto p_t(X_1 = x_1 \mid A_1 = a)\, p_t(A_1 = a).
\end{equation}
$p_t(X_1 = x_1\mid A_1 = a) \propto \exp(-(x_1 - \beta a)^2/(2\Delta))$
is Gaussian in $a$. The prior $p(A_1)$ is the density of
$A_1 = \alpha A_0 + \varepsilon_1$, i.e.\ a Gaussian-Laplace convolution.
This prior on $A_1$ is exactly of the Normal--Laplace type we analysed
in Chapter~\ref{ch:convolution}, with appropriate parameters:
\begin{equation}
A_1 \sim \Normal(\alpha\mu_0, \alpha^2\sigma_0^2) \ast \Lap(0, b).
\end{equation}
The posterior $A_1\mid X_1 = x_1$ is therefore \emph{not} Gaussian.

\section{Proof sketch of the four-term representation}

\begin{proposition}[Four-term representation of $p_t(x_0,x_1,x_2)$]
\label{prop:K2-four-terms}
There exist explicit functions
$\mu_{\sigma_1, \sigma_2}(x_0, x_1, x_2)$ and
$\Sigma_{\sigma_1, \sigma_2}$ for each of the four sign choices
$(\sigma_1, \sigma_2) \in \{+, -\}^2$, such that
\begin{equation}
p_t(x_0, x_1, x_2) =
\sum_{\sigma_1, \sigma_2 \in \{\pm 1\}}
  C_{\sigma_1, \sigma_2}(x_0, x_1, x_2)\,
  \Phi_{2}\!\bigl(
    \mu_{\sigma_1,\sigma_2}(x);\ 0,\ \Sigma_{\sigma_1,\sigma_2}
  \bigr),
\label{eq:K2-four-terms}
\end{equation}
where $\Phi_2(\cdot; 0, \Sigma)$ denotes the bivariate normal CDF with
mean $0$ and covariance $\Sigma$, and $C_{\sigma_1,\sigma_2}$ is an
explicit exponential prefactor depending on the sign pattern.
\end{proposition}

\begin{proof}[Sketch]
Write the joint density in latent form as
\begin{equation}
p_t(x_0,x_1,x_2) = \int\!\!\int p_t(x_0,x_1,x_2, u_1, u_2)\,
                          \mathrm{d}u_1\,\mathrm{d}u_2,
\end{equation}
where $u_i := \varepsilon_i$ and the integrand is the joint density with
the two innovations explicit. The Laplace densities of $u_1, u_2$
contribute $|u_i|$ terms; splitting each by sign,
\begin{equation}
\Lap(u; 0, b) = \sum_{\sigma\in\{\pm 1\}}
  \frac{1}{2b}\,\mathbf{1}_{\{\sigma u \geq 0\}}\,e^{-\sigma u/b}.
\end{equation}
Over each quadrant $\{\sigma_1 u_1 \geq 0,\, \sigma_2 u_2 \geq 0\}$, the
exponent in the integrand is a quadratic-plus-linear function of
$(u_1, u_2)$. Completing the squares (the $K=2$ analogue of the $K=1$
step of Section~\ref{sec:I-plus}) reduces each quadrant integral to a
bivariate normal CDF $\Phi_2$ with an explicit mean vector
$\mu_{\sigma_1, \sigma_2}(x)$ and covariance $\Sigma_{\sigma_1, \sigma_2}$
-- both depending parametrically on the constants $(\alpha, \beta, \sigma_0,
\mu_0, b, \Delta)$. Collecting terms over the four sign choices
gives~\eqref{eq:K2-four-terms}. $\Box$
\end{proof}

\begin{remark}[Why this is as far as we can go in closed form]
The bivariate normal CDF $\Phi_2$ is not elementary (no closed form in
exponentials or rationals), but it is numerically stable via Genz's
algorithm and is treated as a standard special function. Hence the
$K=2$ density is ``as closed-form as one could reasonably hope'': a
finite sum over sign patterns of elementary prefactors times a standard
bivariate Gaussian CDF. At $K=K$, the analogous expression involves a
sum of $2^{K-1}$ terms, each with a $(K-1)$-variate normal CDF. This
exponential-in-$K$ scaling is the \emph{analytical} obstruction to a
full closed form at large $K$: it does not prevent numerical evaluation,
but it does prevent a single tidy expression.
\end{remark}

\section{Practical numerical route}

For simulations and visualisation, the four-term representation is less
useful than direct 3D Fourier inversion of the characteristic
function~\eqref{eq:K2-fourier-original}:
\begin{equation}
p_t(x_0, x_1, x_2) = \frac{1}{(2\pi)^3}\int\!\!\int\!\!\int
    e^{-i(k_0 x_0 + k_1 x_1 + k_2 x_2)}\,\hatfn p_t(k_0, k_1, k_2)\,
    \mathrm{d}k_0\,\mathrm{d}k_1\,\mathrm{d}k_2.
\label{eq:3D-Fourier}
\end{equation}
The integrand decays as $|k_i|^{-2}$ along the innovation frequencies
and as Gaussian along all others, so a uniform grid truncation is fine
for moderate tolerances. In practice we obtain 3--5\% accuracy with
$N^3 = 56^3$ grid points. The $K=2$ numerical script in the bundle
implements this and cross-checks it against a $2\cdot 10^6$-sample Monte
Carlo estimate.

% ============================================================
\chapter{Conjectural rank structure of the $K{=}2$ curvature field}
\label{ch:K2-conj}
% ============================================================

We do not have a full closed-form analogue of
Theorem~\ref{thm:curv} for $K=2$. But the numerical exploration and
structural analysis suggest the following.

\begin{conjecture}[Rank-$\leq 2$ structure of $H_t^{(K=2)}(x)$]
\label{conj:K2-rank}
For the Laplace/OU chain of Chapter~\ref{ch:K2-fourier}, the curvature
field $H_t(x_0, x_1, x_2) = -\nabla^2 \log p_t(x_0,x_1,x_2)$ satisfies
\begin{equation}
H_t(x) = H_t^{\mathrm{site}} + \kappa_t^{(1)}(r_1)\, M_\rho^{(1)}
        + \kappa_t^{(2)}(r_2)\, M_\rho^{(2)} + \mathcal O(\alpha\Delta/v),
\end{equation}
where
\begin{itemize}[leftmargin=2em]
\item $H_t^{\mathrm{site}}$ is a constant $3\times 3$ Gaussian term;
\item each $M_\rho^{(i)}$ is a rank-1 matrix supported on the
nearest-neighbour pair $(i, i+1)$;
\item $\kappa_t^{(1)}(r_1), \kappa_t^{(2)}(r_2)$ are scalar curvature
functions of two residual coordinates $r_1, r_2$;
\item the residual $\mathcal O(\alpha\Delta/v)$ correction captures the
obstruction cross-coupling $M_{12} = -\alpha\Delta$.
\end{itemize}
\end{conjecture}

\begin{interpretation}[What this conjecture says]
In words: the $K=1$ rank-1 result weakens to rank-$\leq 2$ at $K=2$, as
one might expect from two independent innovations. The low-rank
structure is \emph{local} (each rank-1 piece couples only two neighbours),
which is the Markov-locality signature in real space. The deviation from
exactness is controlled by the single parameter
$c = \alpha\Delta/v$ introduced in Section~\ref{sec:residual-setup}.
If this conjecture holds, it suggests a general pattern:
\emph{at $K = K_*$, the curvature field lives on a rank-$\leq K_*$
manifold of matrices, with rank-1 blocks along the chain, and a
controllable perturbation of order $c^{K_*-1}$.}

Testing this conjecture is on the roadmap for the next phase of the
thesis.
\end{interpretation}


% ============================================================
\part{Interpretation and synthesis}
% ============================================================

% ============================================================
\chapter{What Gaussianity buys, what Markovianity buys}
\label{ch:interpretation}
% ============================================================

The central theoretical goal of the thesis programme is to disentangle
the roles of Gaussianity and of Markovianity in shaping the structure of
the joint score. The Laplace AR(1) benchmark is the simplest context in
which these two ingredients can be unmixed. We now summarise what the
exact results in this compendium say on that question.

\section{A brief restatement of the problem}

In the Gaussian AR(1) benchmark (studied separately in
\cite{structure_sigma_t}), the joint density $p_t(x_0, \ldots, x_{K-1})$
is Gaussian, the score $S(x,t)$ is linear in $x$, and the curvature
$H_t = \Sigma_t^{-1}$ is a single matrix whose tridiagonal-mass
fraction $\rho_{\mathrm{tri}}(t)$ measures the persistence of Markov
locality. There are two sources of structure at play:

\begin{itemize}[leftmargin=2em]
\item \textbf{Gaussianity of the innovation and prior} makes the noisy
law Gaussian, hence the score linear, hence the curvature a constant
matrix. It is the reason $\Sigma_t^{-1}$ alone describes everything.
\item \textbf{Markovianity of the clean chain} determines the
\emph{sparsity pattern} of $\Sigma_0^{-1}$ and, more generally, the
\emph{fill-in rate} of $\Sigma_t^{-1}$ as diffusion progresses.
\end{itemize}

These are conceptually distinct. The first is a statement about the
law; the second is a statement about the conditional independence graph.
In the Gaussian case they combine seamlessly, and it is easy to conflate
them. Breaking Gaussianity (by swapping in a Laplace innovation) forces
them apart.

\section{What Gaussianity alone gave us in $K{=}1$}

In the Gaussian AR(1) benchmark at $K=1$: the clean residual is Gaussian,
the OU channel is Gaussian, so everything is Gaussian. The curvature
$\kappa_t^{(G)} = 1/\tau_G^2$ is a constant. The rank-1 structure~\eqref{eq:H-ray}
holds exactly with $\kappa_t^{(G)}$ in place of $\kappa_t(r)$. In this
case, the shape of the curvature matrix -- i.e.\ the rank-1 structure
$M_\rho = nn^\top$ -- is \emph{entirely dictated by Markovianity}, not by
Gaussianity. It reflects the graph structure of the chain: one bond
between $X_0$ and $X_1$, and the residual direction $n \propto (\rho, -1)$
is perpendicular to the residual ray $r = 0$. \emph{This is exactly the
shape we find in the Laplace case, too:} the rank-1 matrix $M_\rho$ is
unchanged by the Laplace-vs-Gaussian choice.

\begin{interpretation}[Markov locality is what gives the rank-1 structure]
The matrix $M_\rho$ in Theorem~\ref{thm:curv} has rank 1 because the chain
has only one bond at $K=1$. This is a structural claim about the graph,
not about the innovation law. Whether the innovation is Gaussian or
Laplace, the curvature field still has rank-1 local structure; what
changes between the two is the scalar in front ($\kappa_t^{(G)}$ constant
vs.\ $\kappa_t(r)$ nonconstant in the Laplace case).
\end{interpretation}

\section{What Gaussianity alone broke when we swapped it for Laplace}

Moving from Gaussian to Laplace kept the rank-1 \emph{shape} but broke the
\emph{constancy} of the scalar in front:
\begin{itemize}[leftmargin=2em]
\item The curvature $\kappa_t(r)$ is now a function of the residual
coordinate, concentrated at $r=0$ and vanishing in the tails.
\item The score $\psi_t(r)$ saturates at $\pm 1/\btilde$ in the tails
rather than growing linearly.
\item The posterior $A_1\mid X_1 = x_1$ is no longer Gaussian at $K=2$,
which is ultimately the reason $M_{12} \neq 0$ there.
\end{itemize}

Notice that the first two effects are \emph{local in $r$} and do not
affect the chain's Markov graph. The third effect, the $K=2$
obstruction, is a genuine new phenomenon: it says that non-Gaussian
priors do not propagate cleanly through OU noising. The single-step
clean analysis is preserved; the multi-step version is not.

\section{What Markovianity kept intact in $K{=}2$}

Even in the $K=2$ case, where Gaussianity fails us, the Markov structure
\emph{still} survives in three ways:
\begin{itemize}[leftmargin=2em]
\item The joint law is still exactly Markov (Proposition~\ref{prop:markov-survives}):
$X_2 \perp X_0 \mid X_1$.
\item The Gaussian sector $M$ of Theorem~\ref{thm:K2-fourier} has a
strict tridiagonal coupling pattern ($M_{02} = 0$) in the rotated
coordinates: there is no direct coupling between the site frequency
$\ell_0$ and the distant innovation $\ell_2$.
\item The rational (non-Gaussian) factor is \emph{exactly diagonal}:
one factor per innovation, independent of all others.
\end{itemize}

The only surviving non-diagonality is the specific cross-term
$M_{12} = -\alpha\Delta$, coming from the shared OU noise at the shared
vertex $A_1$. Markov locality survives; only the particular kind of
``full factorisation into independent bonds'' fails.

\begin{interpretation}[The refined conclusion]
Gaussianity is what gives clean factorisation into independent bonds at
$K=K$; without it, the bonds are diagonal in Fourier space but
\emph{share OU noise} through the shared vertices of the chain. The
graph-theoretic Markov locality (no direct $\ell_0\ell_2$ coupling)
survives the loss of Gaussianity; what fails is the stronger statistical
independence.
\end{interpretation}

\section{Consequences for the thesis research direction}

These observations directly inform the architecture of a learned score
network for dynamic data with temporally correlated frames:

\begin{enumerate}[leftmargin=2em]
\item \textbf{Nearest-neighbour coupling is structurally correct.}
The exact benchmark shows that the right coupling pattern at all $t$ is
nearest-neighbour (tridiagonal in Fourier, at least for AR(1) dynamics
under OU). A score network whose attention pattern respects this
locality should perform well for sequential data, independently of
whether the innovation is Gaussian.

\item \textbf{The innovation-frequency axis is the right coordinate.}
Writing the score in innovation coordinates $(\ell_0, \ell_1, \ell_2, \ldots)$
rather than the raw $(k_0, k_1, k_2, \ldots)$ isolates non-Gaussianity
to a single axis. A score network that factorises along the residual
axis -- rather than along the raw frame axis -- should track the
structure of the exact benchmark.

\item \textbf{Non-Gaussianity is concentrated at finite $t$.}
The centre-benchmark gap $\delta\kappa_t$ scales as $1/a$ at small $t$
and as $1/a^4$ at large $t$. Non-Gaussian fine-tuning is critical at
early diffusion time; at late time the Gaussian benchmark is extremely
accurate, which justifies Gaussian tails in the score network.

\item \textbf{The obstruction sets a complexity floor.}
Exact closed-form solutions at $K \geq 2$ require a finite sum of $2^{K-1}$
multivariate Gaussian CDFs. This exponential-in-$K$ scaling is the
analytic signature of why an exact score network is not feasible for
realistic sequence lengths; it is also why parametric low-rank
approximations -- such as the conjectural rank-$\leq K$ structure of
Conjecture~\ref{conj:K2-rank} -- are the right next step to investigate.
\end{enumerate}

% ============================================================
\chapter{Reading the heatmaps}
\label{ch:heatmaps}
% ============================================================

We close with a pedagogical re-reading of the $K=1$ heatmaps shown in
the main memo (Figure 2). The goal is to explain, step by step, what
the geometry of the two-dimensional curvature field is telling us.

\section{The residual ray as a one-dimensional manifold}

The residual coordinate $r = x_1 - \nu - \rho x_0$ vanishes on the line
$x_1 = \nu + \rho x_0$ in the $(x_0, x_1)$ plane. This line is the
\emph{residual ray} -- a one-dimensional manifold through the bivariate
domain on which the Laplace cusp lives.

The slope of the ray is $\rho$, which interpolates between $\alpha$ at
$t=0$ (the clean AR(1) slope) and $0$ as $t\to\infty$ (when the noisy
$X_1$ decouples from $X_0$).

\section{Interpretation of $H_{11}(x)$}

The entry $H_{11}(x) = \kappa_t(r)$ depends on $(x_0, x_1)$ only through
the residual $r$. So $H_{11}$ is constant along \emph{any} line parallel
to the residual ray, i.e.\ along the family of lines
$x_1 = \nu + \rho x_0 + r$. Geometrically, $H_{11}$ shows a
``stripe'' pattern: high values concentrated in a band around the
residual ray, decaying to zero as one moves away.

\begin{whybox}
Why stripes? Because the non-Gaussianity lives entirely in the one
residual direction. Anywhere you can move \emph{parallel to} the ray,
the residual $r$ stays the same and so does $\kappa_t(r)$. You must move
\emph{perpendicular to} the ray to see $\kappa_t(r)$ change.
\end{whybox}

\section{Interpretation of $H_{01}(x)$}

The off-diagonal entry is
$H_{01}(x) = -\rho \kappa_t(r)$. It is a \emph{multiple} of $H_{11}$,
with the same stripe pattern and proportionality $-\rho$. A positive
$\rho$ (e.g.\ $\alpha > 0$) gives a negative off-diagonal; a negative
$\rho$ gives a positive off-diagonal; and $\rho = 0$ gives no
off-diagonal at all. At $t=0$ where $\rho = \alpha$, the off-diagonal
reflects the clean correlation. At $t\to\infty$ where $\rho\to 0$, the
off-diagonal vanishes -- the two coordinates decouple.

\section{What would change if we swapped the innovation law?}

Under the exchange Laplace $\to$ Gaussian: the heatmaps become
\emph{constant} in color, with the stripe pattern replaced by a single
uniform value. The rank-1 structure is unchanged; only the scalar
$\kappa_t$ in front of it flips from a $r$-dependent Laplace curvature
to a constant Gaussian curvature. This is the graphical embodiment of
the interpretation of Section~\ref{ch:interpretation}: Markov locality
$\to$ stripes; Gaussianity $\to$ the color of each stripe is constant.

Under a general non-Gaussian innovation: the stripe pattern persists
(Markov still holds), but the \emph{shape} of $\kappa_t(r)$ as a function
of $r$ changes. For multimodal innovations, $\kappa_t$ can even be
negative, indicating regions of the residual-ray direction where the
density is locally convex in $\log$-space (``anti-confining'' regions).
Investigating this is listed as one of the open questions in the main
memo.

% ============================================================
\chapter{Open questions revisited}
\label{ch:open}
% ============================================================

The main memo lists five open questions. Here we indicate, for each, what
specific missing piece of the exact picture it corresponds to and how
it might be attacked.

\begin{enumerate}[leftmargin=2em]
\item \textbf{Centre-tail gap scaling for general innovation laws.}
From the ODE~\eqref{eq:h-ODE-repeat}, $h_t''(0) = (h_t(0) - \phi_\tau(0))/\btilde^2$.
Thus $\kappa_t(0) = -h_t''(0)/h_t(0) = -(1 - \phi_\tau(0)/h_t(0))/\btilde^2$.
This is a universal formula: for \emph{any} innovation law with a density
that we can evaluate at $0$, the centre curvature is determined by
$h_t(0)$. So the centre-tail gap $\kappa_t(0) - \kappa_t^{(G)}$ reduces to
studying the functional $\varepsilon \mapsto h_t(0)$. For a variance-matched
Gaussian, this ratio is always $\sqrt{2/\pi}/a$ at leading order at small
$a$, with model-dependent subleading corrections controlled by higher
cumulants. A rigorous treatment is accessible via Edgeworth expansion.

\item \textbf{Variance-Gamma bridge.}
The VG family $\hatfn f_\varepsilon(\ell) = (1 + b^2\ell^2/(2\nu))^{-\nu}$
interpolates Laplace ($\nu=1$) to Gaussian ($\nu\to\infty$). At integer
$\nu$, the residual density $h_t^{(\nu)}$ is a sum of $\nu$ Laplace-type
terms (via partial fractions), and the centre curvature admits a closed
form as a generalisation of~\eqref{eq:mills-main}. Computing this and
tracking $\kappa_t^{(\nu)}(0)$ along the VG family would quantify how
the rank-1 scalar $\kappa_t$ turns on with non-Gaussianity.

\item \textbf{Multimodal innovations.}
For the symmetric Gaussian-mixture innovation
$\varepsilon \sim \tfrac12 \Normal(-m, \sigma^2) + \tfrac12 \Normal(+m, \sigma^2)$,
the characteristic function is
$\hatfn f_\varepsilon(\ell) = \cos(m\ell)e^{-\sigma^2\ell^2/2}$, and the
residual $R_t = \beta\varepsilon + G_t$ has density
\begin{equation}
h_t(r) = \frac{1}{2}[\phi_{s_t}(r - \beta m) + \phi_{s_t}(r + \beta m)],
\qquad s_t^2 = \tau^2 + \beta^2 \sigma^2.
\end{equation}
For small $t$, $h_t$ has two peaks at $\pm\beta m$; for large $t$, they
merge. $\kappa_t(r)$ can be \emph{negative} in a range of $r$. The
critical diffusion time $t^*$ at which the bimodality disappears is an
explicit function of $m/\sigma$.

\item \textbf{Nonlinear locality observables.}
Define
$T_d(x, t) := \sum_{|i-j|\leq d} H_t(x)_{ij}^2 / \|H_t(x)\|_F^2$,
the fractional Frobenius mass supported within a band of width $d$.
In the Gaussian benchmark, $T_d(x, t)$ is constant in $x$ and matches
the Gaussian tridiagonal-mass measure $\rho_{\mathrm{tri}}(t)$. In the
Laplace benchmark, $T_d(x,t)$ depends on $x$ through $\kappa_t(r)$; its
pointwise expectation is
$\bar T_d(t) = \E_{X\sim p_t}[T_d(X,t)]$, a true functional replacement
for the Gaussian measure. Computing $\bar T_d$ as a function of $t$ and
$d$ is the bridge between the Gaussian benchmark and the non-Gaussian
benchmark.

\item \textbf{Testing Conjecture~\ref{conj:K2-rank}.}
Systematic evaluation of $H_t$ on a 3D grid of $(x_0, x_1, x_2)$ values
(say $30^3 = 27000$ points), SVD of each $3\times 3$ matrix, statistical
analysis of the resulting rank distributions. The conjecture predicts
that the rank-$\leq 2$ approximation error decays as $(\alpha\Delta/v)^2$.
Verifying this numerically is a direct test.
\end{enumerate}

\begin{thebibliography}{9}
\bibitem{laplace_note}
G.~Mantovani, M.~Lomel\'e,
\emph{Laplace AR(1) under Ornstein--Uhlenbeck Corruption: Detailed Derivation Note},
internal manuscript, 2026.

\bibitem{short_presentation}
G.~Mantovani, M.~Lomel\'e,
\emph{Markov Locality, Diffusion Fill-In, and the First Non-Gaussian Route},
short presentation note, March 2026.

\bibitem{structure_sigma_t}
G.~Mantovani,
\emph{Structure of $\Sigma_t$ and $\Sigma_t^{-1}$ in Gaussian AR(1) Diffusion},
internal note, 2026.

\bibitem{thesis_part1}
G.~Mantovani, M.~Lomel\'e,
\emph{Score-Based Generative Diffusion Models for Dynamic Objects},
Thesis Part I, Bocconi University, 2026.
\end{thebibliography}

\end{document}
